\chapter{Algebraic arguments, integrals and ladders}
\label{ch:03-algebraic}

Away from roots of unity, functional equations still move values onto
structured sets of arguments. Vertical lines and half-angle substitutions give
tractable examples; unit relations in number fields lead to longer ladders.
The organizing principle is to understand the map of arguments before choosing
the constants in an integer-relation search.

\section{The problem and the key geometric fact}
\href{https://math.stackexchange.com/q/1410276}{MSE 1410276} records
\begin{equation}\label{vertical:eq:mse}
  \Real\Li_2(1\pm\iu\sqrt3)=\frac{\pi^2}{24}-\tfrac12\ln^2 2-\tfrac14\Li_2(\tfrac14),
\end{equation}
with $\Imag\Li_2(1+\iu\sqrt3)=\tfrac{\pi\ln2}{2}+\tfrac{5\sqrt3}{72}\bigl[\psi_1(\tfrac13)-\psi_1(\tfrac23)\bigr]$
(\href{https://math.stackexchange.com/a/1410829}{a/1410829}, \href{https://math.stackexchange.com/a/1410705}{a/1410705}).
The point $1+\iu\sqrt3=2e^{\iu\pi/3}$, but the structurally decisive fact is that for any $z=1+\iu y$ the
reflection point $1-z=-\iu y$ is \emph{purely imaginary}. This is exactly what makes the dilogarithm
functional equations collapse.

\section{The weight-2 vertical-line theorem}
\begin{theorem}\label{vertical:thm:wt2}
For $y>0$,
\begin{align}
  \Real\Li_2(1+\iu y)&=\frac{\pi^2}{6}-\tfrac12\ln(1{+}y^2)\ln y-\frac\pi2\arctan y-\tfrac14\Li_2(-y^2),
  \label{vertical:eq:Re2}\\
  \Imag\Li_2(1+\iu y)&=\frac\pi4\ln(1{+}y^2)-\arctan y\,\ln y+\Ti_2(y),\qquad \Ti_2(y)=\Imag\Li_2(\iu y).
  \label{vertical:eq:Im2}
\end{align}
\end{theorem}
\begin{proof}
The reflection law $\Li_2(z)+\Li_2(1-z)=\zeta(2)-\ln z\ln(1-z)$ at $z=1+\iu y$, $1-z=-\iu y$, gives
$\Li_2(1+\iu y)=\zeta(2)-\ln(1+\iu y)\ln(-\iu y)-\Li_2(-\iu y)$. The duplication law
$\Li_2(\iu y)+\Li_2(-\iu y)=\tfrac12\Li_2(-y^2)$ together with $\overline{\Li_2(w)}=\Li_2(\bar w)$ yields
$\Real\Li_2(-\iu y)=\tfrac12\bigl(\Li_2(\iu y)+\Li_2(-\iu y)\bigr)=\tfrac14\Li_2(-y^2)$ and
$\Imag\Li_2(-\iu y)=-\Ti_2(y)$. Expanding $\ln(1+\iu y)=\tfrac12\ln(1{+}y^2)+\iu\arctan y$ and
$\ln(-\iu y)=\ln y-\tfrac{\iu\pi}2$ and taking real and imaginary parts gives~\eqref{vertical:eq:Re2}--\eqref{vertical:eq:Im2}.
\end{proof}
On the line $\Real z=1$ the dilogarithm thus splits into a \emph{real} part carried by the even real
dilogarithm $\Li_2(-y^2)$ and an \emph{imaginary} part carried by the odd inverse-tangent integral
$\Ti_2(y)$. Equation~\eqref{vertical:eq:mse} is the case $y=\sqrt3$ (using
$\Li_2(-3)=\Li_2(\tfrac14)-\zeta(2)+2\ln2\ln\tfrac23$), and the trigamma form of the imaginary part is just
$\Ti_2(\sqrt3)=\tfrac{\pi}6\ln3+\tfrac{5\sqrt3}{72}[\psi_1(\tfrac13)-\psi_1(\tfrac23)]$ substituted
into~\eqref{vertical:eq:Im2}.

\section{The Clausen form at $y=\tan\theta$}
Since $1+\iu\tan\theta=e^{\iu\theta}/\cos\theta$, Theorem~\ref{vertical:thm:wt2} combines with Ramanujan's
$\Ti_2(\tan\theta)=\theta\ln\tan\theta+\tfrac12\bigl[\Cl_2(2\theta)+\Cl_2(\pi-2\theta)\bigr]$ so that the
$\theta\ln\tan\theta$ terms cancel, leaving a pure-Clausen imaginary part:
\begin{equation}\label{vertical:eq:clausen}
  \Imag\Li_2(1+\iu\tan\theta)=-\frac\pi2\ln\cos\theta+\tfrac12\bigl[\Cl_2(2\theta)+\Cl_2(\pi-2\theta)\bigr].
\end{equation}
\begin{corollary}\label{vertical:cor:specials}
$\Real\Li_2(1+\iu)=\dfrac{\pi^2}{16}$ and $\Imag\Li_2(1+\iu)=\dfrac{\pi\ln2}{4}+G$ (with $G$ Catalan's
constant), from $y=1$ in Theorem~\ref{vertical:thm:wt2} (using $\Li_2(-1)=-\pi^2/12$, $\Ti_2(1)=G$).
\end{corollary}

\section{Weight 3: a half-lifted symmetry}
At weight $3$ the real/imaginary split survives only in part. The \emph{real} part still closes at special
$y$:
\begin{equation}\label{vertical:eq:Re3}
  \Real\Li_3(1+\iu)=\frac{35}{64}\zeta(3)+\frac{\pi^2\ln2}{32},\qquad
  \Real\Li_3(1+\iu\sqrt3)=\frac{35\zeta(3)+18\Li_3(\tfrac14)+6\pi^2\ln2-24\ln^3 2}{144},
\end{equation}
the second being the exact weight-3 analogue of the $\Li_2(\tfrac14)$ in~\eqref{vertical:eq:mse}. For the general
real part, the displayed calculation introduces a further integral: differentiating, $\partial_y\Real\Li_3(1+\iu y)=-\bigl(\Imag\Li_2(1+\iu y)-y\,\Real\Li_2(1+\iu y)\bigr)/(1+y^2)$,
and integrating the $\Ti_2$ contribution produces $\int_0^y\arctan^2(t)/t\,dt$, a weight-three inverse-tangent-square integral. The displayed special cases at $y=1,\sqrt3$ simplify, without giving an elementary-value classification for every rational angle. The \emph{imaginary} part at $1+\iu$ is represented here by the
central-binomial Gaussian tower constant
\begin{equation}\label{vertical:eq:Im3}
  \Imag\Li_3(1+\iu)=\mathcal F_3-\frac{\pi^3}{192},\qquad
  \mathcal F_3=\sqrt2\,{}_4F_3\bigl(\{\tfrac12\}^4;\{\tfrac32\}^3;\tfrac12\bigr),
\end{equation}
the $n=3$ entry of the $\Imag\Li_n(1+\iu)$ tower (companion report
\texttt{gaussian-1plusi-tower}). So on $\Real z=1$ the imaginary part is precisely the deep object that
the central-binomial ladder measures, while the real part remains tractable---the two halves of the
polylogarithm part ways at weight $3$.
\begin{remark}[An unsupported nonreduction claim]\label{vertical:rem:nonreduction}
Earlier prose asserted that this imaginary part was not elementary, without
specifying a comparison algebra or supplying an obstruction. That assertion
is an open reducibility question, rather than a consequence of
\eqref{vertical:eq:Im3}. The displayed hypergeometric representation remains
valid; any proposed reduction or nonreduction must name its algebra of
constants and establish a separate proof.
\end{remark}
\section{The integral and its provenance}
The object (Math.StackExchange \href{https://math.stackexchange.com/q/564816}{question 564816}) is
\begin{equation}\label{cleo:eq:I}
  I=\int_0^{\pi/2}\arctan^2\!\left(\frac{6\sin x}{3+\cos2x}\right)dx=1.330970396842566666125327\ldots
\end{equation}
\href{https://math.stackexchange.com/a/570815}{Cleo} gave the dilogarithmic form
\begin{multline}\label{cleo:eq:cleo}
  I=\frac{\pi}{2}\Bigl(2\Li_2\bigl((1-\sqrt2)(2-\sqrt5)\bigr)-2\Li_2\bigl((1-\sqrt2)(\sqrt5-2)\bigr)\\
  +\Li_2(3-\sqrt8)-\Li_2(\sqrt8-3)+\Li_2(9-\sqrt{80})-\Li_2(\sqrt{80}-9)\Bigr),
\end{multline}
while \href{https://math.stackexchange.com/a/571079}{Sangchul Lee} gave the Legendre-chi form
\begin{equation}\label{cleo:eq:lee}
  I=\pi\Bigl(\chitwo(3-2\sqrt2)+\chitwo(9-4\sqrt5)+2\chitwo\bigl((\sqrt2-1)(\sqrt5-2)\bigr)\Bigr),
\end{equation}
where $\chitwo(x)=\tfrac12\bigl(\Li_2(x)-\Li_2(-x)\bigr)=\sum_{k\ge0}\frac{x^{2k+1}}{(2k+1)^2}$. These are
identical: $\Li_2(x)-\Li_2(-x)=2\chitwo(x)$ turns each $\Li_2$ pair in~\eqref{cleo:eq:cleo} into a $\chitwo$, and
$3-\sqrt8=(\sqrt2-1)^2$, $9-\sqrt{80}=(\sqrt5-2)^2$, $(1-\sqrt2)(2-\sqrt5)=(\sqrt2-1)(\sqrt5-2)$, so both
equal
\begin{equation}\label{cleo:eq:canon}
  I=\pi\Bigl[\chitwo\bigl((\sqrt2-1)^2\bigr)+\chitwo\bigl((\sqrt5-2)^2\bigr)
  +2\,\chitwo\bigl((\sqrt2-1)(\sqrt5-2)\bigr)\Bigr].
\end{equation}
The arguments are a silver$\,\otimes\,$golden pair: $\sqrt2-1=\tan\frac\pi8=\dS^{-1}$ (reciprocal silver
ratio $\dS=1+\sqrt2$) and $\sqrt5-2=\phi^{-3}$ (golden $\phi=\tfrac{1+\sqrt5}2$, $\phi^3=\sqrt5+2$), so
$I=\pi[\chitwo(\dS^{-2})+\chitwo(\phi^{-6})+2\chitwo(\dS^{-1}\phi^{-3})]$.

\section{The master formula (Reshetnikov--Lee) and a differentiation proof}
The closed forms~\eqref{cleo:eq:cleo}--\eqref{cleo:eq:lee} are instances of a two-parameter master formula. Its
seed is Reshetnikov's single-arctangent ladder
$\int_0^{\pi/2}\arctan(t\sin x)\,dx=2\chitwo(u_t)$; inspired by it, Lee
(\href{https://math.stackexchange.com/a/562492}{a/562492}, 2013) conjectured and proved
\begin{theorem}[Lee 2013]\label{cleo:thm:master}
For real $p,q$,
\begin{equation}\label{cleo:eq:master}
\begin{aligned}
  J(p,q)&:=\int_0^{\pi/2}\arctan(p\sin x)\,\arctan(q\sin x)\,dx=\pi\,\chitwo\!\bigl(\up\uq\bigr),\\
  &\qquad\text{where}\quad\up=\frac{\sqrt{1+p^2}-1}{p}=\tan\!\bigl(\tfrac12\arctan p\bigr).
\end{aligned}
\end{equation}
\end{theorem}
Lee gives two proofs (a triple integral and a Fourier-orthogonality argument). We record a short third
proof by differentiation, which is the most direct route and is machine-verifiable.
\begin{proof}
Both sides vanish on the axes, and $\partial_pJ(p,0)=\partial_p[\pi\chitwo(\up u_0)]=0$ (since $u_0=0$),
symmetrically in $q$; hence each side equals the double integral of its mixed partial over
$[0,p]\times[0,q]$, and it suffices to match mixed partials. Differentiating under the integral and using
the partial fraction $\frac{s^2}{(1+p^2s^2)(1+q^2s^2)}=\frac{1}{p^2-q^2}\bigl(\frac1{1+q^2s^2}-\frac1{1+p^2s^2}\bigr)$
together with $\int_0^{\pi/2}\frac{dx}{1+c\sin^2x}=\frac{\pi}{2\sqrt{1+c}}$,
\begin{equation}\label{cleo:eq:lhsmixed}
  \partial_p\partial_q J=\frac{\pi}{2(p^2-q^2)}\Bigl(\frac{1}{\sqrt{1+q^2}}-\frac{1}{\sqrt{1+p^2}}\Bigr).
\end{equation}
On the right, $\chitwo'(w)=\arctanh(w)/w$ gives
$\partial_p\partial_q[\pi\chitwo(\up\uq)]=\pi\,\up'\uq'/(1-\up^2\uq^2)$. The substitution
$p=\frac{2\up}{1-\up^2}$ yields $\sqrt{1+p^2}=\frac{1+\up^2}{1-\up^2}$ and $\up'=\frac{(1-\up^2)^2}{2(1+\up^2)}$,
whence both sides reduce to the same rational function
$\frac{\pi(1-\up^2)^2(1-\uq^2)^2}{4(1+\up^2)(1+\uq^2)(1-\up^2\uq^2)}$. (Wolfram confirms this reconciliation
symbolically: \texttt{FullSimplify} of the difference returns $0$.)
\end{proof}

\section{Cleo's integral as the perfect-square instance}\label{cleo:sec:square}
The argument in~\eqref{cleo:eq:I} collapses to a sum of arctangents. With $\cos2x=1-2\sin^2x$,
\[
  \frac{6\sin x}{3+\cos2x}=\frac{3\sin x}{2-\sin^2x}=\frac{(1+\tfrac12)\sin x}{1-(1\cdot\tfrac12)\sin^2x},
\]
the tangent-addition shape with roots $1,\tfrac12$ of $t^2-\tfrac32t+\tfrac12$, so for $0\le x\le\frac\pi2$
\begin{equation}\label{cleo:eq:collapse}
  \arctan\!\left(\frac{6\sin x}{3+\cos2x}\right)=\arctan(\sin x)+\arctan(\tfrac12\sin x),
\end{equation}
so that $I=\int_0^{\pi/2}\bigl(\arctan(\sin x)+\arctan(\tfrac12\sin x)\bigr)^2 dx$.
Applying Theorem~\ref{cleo:thm:master} to the expanded square, with $u_1=\sqrt2-1$, $u_{1/2}=\sqrt5-2$,
reproduces~\eqref{cleo:eq:canon}. The coefficient pattern $1,2,1$ is a perfect square: writing $a=\sqrt2-1$,
$b=\sqrt5-2$,
\begin{equation}\label{cleo:eq:squareseries}
  \frac{I}{\pi}=\sum_{k\ge0}\frac{\bigl((\sqrt2-1)^{2k+1}+(\sqrt5-2)^{2k+1}\bigr)^2}{(2k+1)^2}.
\end{equation}
The coefficients $a^{2k+1}+b^{2k+1}$ are the $\sin$-power Fourier coefficients of
$\arctan(\sin x)+\arctan(\tfrac12\sin x)$; \eqref{cleo:eq:squareseries} is the Parseval form underlying Lee's
Fourier proof.

\section{Closed values and a bounded experimental survey}
Three classical values supply useful exact specializations. With
$\phi=(1+\sqrt5)/2$,
\begin{equation}\label{cleo:eq:elem}
\begin{aligned}
\chi_2(\phi^{-1})&=\frac{\pi^2}{12}-\frac34\log^2\phi,\\
\chi_2(\phi^{-3})&=\frac{\pi^2}{24}-\frac34\log^2\phi,\\
\chi_2(\sqrt2-1)&=\frac{\pi^2}{16}-\frac14\log^2(1+\sqrt2).
\end{aligned}
\end{equation}
The source surveyed inverse metallic means $m_n^{-k}$ with $1\le n,k\le6$,
selected rational points, and $\tan(\pi/k)$ for $k=5,8,10,12$ and their
squares. Searches at $130$--$140$ digits found no further accepted
reduction in their selected logarithmic and $\pi^2$ baskets. This does not
classify all elementary values. Nor would independence of the individual
terms alone prove that their sum cannot simplify.

\section{New closed-form integrals from elementary $\chitwo$}
Conversely,~\eqref{cleo:eq:elem} feeds the master formula to produce integrals that \emph{do} close. The
cleanest cross integral uses $u_2=\phi^{-1}$, $u_{2/\sqrt5}=\phi^{-2}$, so $u_2u_{2/\sqrt5}=\phi^{-3}$:
\begin{equation}\label{cleo:eq:A}
  \boxed{\;\int_0^{\pi/2}\arctan(2\sin x)\,\arctan\!\bigl(\tfrac2{\sqrt5}\sin x\bigr)\,dx
  =\frac{\pi^3}{24}-\frac{3\pi}{4}\ln^2\phi\;}\qquad(=0.7463164407\ldots).
\end{equation}
The diagonal $\int_0^{\pi/2}\arctan^2(\lambda\sin x)\,dx=\pi\,\chitwo(u_\lambda^2)$ closes at three clean
$\lambda$, since $u_{\sqrt\phi}^2=\phi^{-3}$, $u_{\sqrt{2+2\sqrt2}}^2=\sqrt2-1$, and $u_{2\phi^{3/2}}^2=\phi^{-1}$:
\begin{equation}\label{cleo:eq:diag}
  \int_0^{\pi/2}\!\arctan^2(\sqrt\phi\,\sin x)\,dx=\frac{\pi^3}{24}-\frac{3\pi}{4}\ln^2\phi,\quad
  \int_0^{\pi/2}\!\arctan^2\!\bigl(2\phi^{3/2}\sin x\bigr)\,dx=\frac{\pi^3}{12}-\frac{3\pi}{4}\ln^2\phi.
\end{equation}
The simplest non-elementary diagonal is the $(1,1)$ case $\int_0^{\pi/2}\arctan^2(\sin x)\,dx=\pi\,\chitwo(3-2\sqrt2)$.

\section{A sibling trichotomy}
The master formula has two siblings under the companion half-angle substitution
$\vp=\frac{1-\sqrt{1-p^2}}{p}=\tanh(\tfrac12\arctanh p)$ (for $|p|<1$). Replacing one or both
arctangents by $\arctanh$ produces:
\begin{equation}\label{cleo:eq:trich}
  \int_0^{\pi/2}\!\!\begin{cases}\arctan(p\sin x)\,\arctan(q\sin x)\\[2pt]
  \arctanh(p\sin x)\,\arctanh(q\sin x)\\[2pt]
  \arctan(p\sin x)\,\arctanh(q\sin x)\end{cases}\!\!dx=
  \begin{cases}\pi\,\chitwo(\up\uq),\\[2pt]\pi\,\chitwo(\vp\vq),\\[2pt]\pi\,\Ti_2(\up\vq),\end{cases}
\end{equation}
where $\Ti_2(z)=\Im\Li_2(\mathrm iz)=\int_0^z\frac{\arctan t}{t}\,dt$ is the inverse-tangent integral: like
pairs give the Legendre chi, the mixed pair gives the inverse-tangent integral. The linear shadows are
$\int_0^{\pi/2}\arctan(\lambda\sin x)\,dx=2\chitwo(u_\lambda)$ (Reshetnikov) and, with $\theta=\arcsin\lambda$,
\begin{equation}\label{cleo:eq:lintanh}
  \int_0^{\pi/2}\arctanh(\lambda\sin x)\,dx=\theta\,\ln\tan\tfrac\theta2+2\Cl_2(\theta)-\tfrac12\Cl_2(2\theta),
\end{equation}
the latter showing that the $\arctanh$ branch is \emph{not} a clean $2\chitwo(v_\lambda)$ mirror---that
coincidence is special to $\arctan$. By $x\mapsto\frac\pi2-x$ the $\cos$ versions of all of these equal the
$\sin$ versions. The weight-$3$ analogue $\int_0^{\pi/2}\arctan(p\sin x)\arctan(q\sin x)/\sin x\,dx$ does not factor by the same mechanism into a single $\chi_3$: its coefficient $\frac{4^N}{(2N+1)\binom{2N}{N}}$ couples the two
summation indices only through $N=m+n$, so it does not factor through $\up\uq$ (numerically no clean
$\chi_3$ relation appears up to height $10^8$).

\section{An alternate form, and minimality}\label{cleo:sec:minimal}
\paragraph{Landen-mirror form.} Applying the $\chitwo$ Landen law
\begin{equation}\label{cleo:eq:landen}
  \chitwo(x)+\chitwo\!\left(\frac{1-x}{1+x}\right)=\frac{\pi^2}{8}-\frac12\ln x\,\ln\frac{1-x}{1+x}
  \qquad(0<x<1)
\end{equation}
to the three arguments of~\eqref{cleo:eq:canon} sends $(\sqrt2-1)^2\mapsto\tfrac1{\sqrt2}$,
$(\sqrt2-1)(\sqrt5-2)\mapsto\sqrt5-\sqrt2$, $(\sqrt5-2)^2\mapsto\tfrac2{\sqrt5}$, giving a second three-term
representation
\begin{equation}\label{cleo:eq:mirror}
  \frac{I}{\pi}=\mathcal E-\Bigl[\chitwo(\tfrac1{\sqrt2})+2\,\chitwo(\sqrt5-\sqrt2)+\chitwo(\tfrac2{\sqrt5})\Bigr],
\end{equation}
with the elementary constant
$\mathcal E=\frac{\pi^2}{2}-\frac12\ln2\,\ln(1+\sqrt2)+3\ln2\,\ln\phi-\frac32\ln5\,\ln\phi
+\bigl(\ln(1+\sqrt2)+3\ln\phi\bigr)\ln(\sqrt5-\sqrt2)$.
This representation trades the small original arguments for larger ones
and explicit logarithms. It follows directly by applying the Landen law
three times.

\paragraph{Search for shorter forms.}
The source tested $209$ algebraic arguments in
$\mathbb Q(\sqrt2,\sqrt5,\sqrt{10})$ at at least $120$ digits, with
residual threshold $10^{-105}$. It found no accepted one-term or two-term
representation in the tested baskets. The three-term formula is an exact
evaluation, while its minimality and non-elementarity remain unproved.

\section{Unit relations and the golden family}
With $\phi=\tfrac{1+\sqrt5}2$, $\rho=\phi^{-1}$, the defining relation $1-\rho^2=\rho$ is the seed of the
golden ladders. Reshetnikov's conjectures live among the values $\Li_n(\phi^{-2k})$.
The source used PSLQ to recover coefficient vectors and tested the resulting
relations numerically (\texttt{mpmath}, $130$--$420$ digits, Champernowne
canary), with Wolfram comparisons where the cancellation depth allowed.
Exact functional-equation certificates, where supplied, are identified
separately. This description concerns the evidence presented here; it does
not assert that the corresponding identities lack proofs in the literature.
\begin{remark}
The source reports cancellation of $\sim10^8$-sized terms with more
than $100$ matching decimal digits. Its historical Wolfram evaluation
under-resolved some cancellations, so independently controlled precision
in mpmath supplied a comparator. Agreement between numerical engines is
a useful check, but does not establish exact equality.
\end{remark}

\section{Reshetnikov's formulas and their numerical evidence}
The source recovered the following tetralogarithm relation
(\href{https://math.stackexchange.com/q/1373123}{MSE 1373123},
\href{https://mathoverflow.net/q/261408}{MO 261408}) numerically.
The exact proof below now establishes it. With $\alpha=\ln2,\beta=\ln3$, its form is
\begin{multline}\label{golden:eq:tetra}
  720\Li_4(\tfrac12)-2160\Li_4(\tfrac13)+2160\Li_4(\tfrac23)\\+270\Li_4(\tfrac14)+540\Li_4(\tfrac34)+135\Li_4(\tfrac19)\\
  =19\pi^4+30(\pi^2-\beta^2)(10\alpha^2-12\alpha\beta+3\beta^2)-30\alpha^2(19\alpha^2-24\alpha\beta+8\beta^2),
\end{multline}
The source recovered this relation in its weight-four basket; Theorem~\ref{rattetra:thm:main} proves it without a numerical-recognition premise. Uniqueness of all relations among the individual values remains a separate question.
For the three $\Li_5$ golden formulas below, the source reports residuals
of approximately $10^{-257}$. Their connection with the ladders of
Abouzahra and Lewin~\cite{golden:AbouzahraLewin1985} is described below:
{\scriptsize
\begin{align}
  4455\Li_5(\phi^{-2})-1215\Li_5(\phi^{-4})-360\Li_5(\phi^{-6})+15\Li_5(\phi^{-12})
  &=3015\zeta(5)-702\lp^5+180\pi^2\lp^3-38\pi^4\lp,\notag\\
  45000\Li_5(\phi^{-2})-16875\Li_5(\phi^{-4})-144\Li_5(\phi^{-10})+9\Li_5(\phi^{-20})
  &=28944\zeta(5)-6000\lp^5+1600\pi^2\lp^3-356\pi^4\lp,\notag\\
  15660\Li_5(\phi^{-2})-19440\Li_5(\phi^{-4})+7680\Li_5(\phi^{-6})
  &\notag\\[-2pt]
  \quad+2430\Li_5(\phi^{-8})-15\Li_5(\phi^{-24})
  &=5025\zeta(5)+2088\lp^5-240\pi^2\lp^3-28\pi^4\lp.\notag
\end{align}}

\section{Reconstructing Tito's weight-6 and weight-7 formulas}
Tito's ``old answer'' lists the golden family at weights $2$--$5$
(numerically checked in the source), then states the weight-$6$ and
weight-$7$ ladders only schematically, the integer coefficients ``too
tedious to write.'' PSLQ recovers the following candidate coefficient
vectors on the six-argument set $\{\phi^{-2},\phi^{-4},\phi^{-6},\phi^{-8},\phi^{-12},\phi^{-24}\}$:
{\footnotesize
\begin{align}
\label{golden:eq:Li6}
  &98820\Li_6(\phi^{-2})-104490\Li_6(\phi^{-4})+33120\Li_6(\phi^{-6})+7290\Li_6(\phi^{-8})-50\Li_6(\phi^{-12})-15\Li_6(\phi^{-24})\notag\\
  &\qquad=-13032\lp^6+3240\pi^2\lp^4-424\pi^4\lp^2+\mathbf{32155}\,\zeta(6),\\
\label{golden:eq:Li7}
  &74707920\Li_7(\phi^{-2})-39497220\Li_7(\phi^{-4})+8346240\Li_7(\phi^{-6})+1377810\Li_7(\phi^{-8})-6300\Li_7(\phi^{-12})-945\Li_7(\phi^{-24})\notag\\
  &\qquad=2814912\lp^7-979776\pi^2\lp^5+213696\pi^4\lp^3-51448\pi^6\lp+\mathbf{44689995}\,\zeta(7).
\end{align}}
The $\zeta(6)$ and $\zeta(7)$ coefficients reproduce Tito's stated
$32155$ and $44689995$ exactly. The source also reports direct-substitution
residuals $10^{-117}$ and $10^{-114}$ for
\eqref{golden:eq:Li6}--\eqref{golden:eq:Li7}, with primitive coefficient
vectors. These are consistent numerical checks; an exact derivation of
these displayed reconstructions is not supplied in this section.

\section{Explicit candidate formulas at weights 8 and 9}
Tito remarks that the index-$24$ ladder ``can be extended even to $n=9$''
but gives no formula. The source recovers the following candidates.
Its searches on the six-argument set of
\eqref{golden:eq:Li6}--\eqref{golden:eq:Li7} detected numerical reductions
at weights $2$--$7$ and none at $8$ or $9$ within the tested bounds.
The accepted higher-weight candidates use the \emph{eight}-argument union
of the index-$12$, $20$, $24$
sets, $\mathcal A=\{\phi^{-2},\phi^{-4},\phi^{-6},\phi^{-8},\phi^{-10},\phi^{-12},\phi^{-20},\phi^{-24}\}$:
{\scriptsize
\begin{align}
\label{golden:eq:Li8}
  &{-}490795200\Li_8(\phi^{-2})+636131475\Li_8(\phi^{-4})-388684800\Li_8(\phi^{-6})+82668600\Li_8(\phi^{-8})\notag\\
  &\quad+9507456\Li_8(\phi^{-10})+5250000\Li_8(\phi^{-12})-74277\Li_8(\phi^{-20})-18900\Li_8(\phi^{-24})\notag\\
  &\qquad=-23641920\lp^8+7956480\pi^2\lp^6-1434720\pi^4\lp^4+184432\pi^6\lp^2-148565886\,\zeta(8),\\
\label{golden:eq:Li9}
  &220857840000\Li_9(\phi^{-2})-143129581875\Li_9(\phi^{-4})+58302720000\Li_9(\phi^{-6})-9300217500\Li_9(\phi^{-8})\notag\\
  &\quad-855671040\Li_9(\phi^{-10})-393750000\Li_9(\phi^{-12})+3342465\Li_9(\phi^{-20})+708750\Li_9(\phi^{-24})\notag\\
  &\qquad=-2364192000\lp^9+1022976000\pi^2\lp^7-258249600\pi^4\lp^5\notag\\ &\qquad\quad+55329600\pi^6\lp^3-14149132\pi^8\lp+125596001790\,\zeta(9).
\end{align}}
The source reports a residual printed as $0$ at $150$ working digits for
\eqref{golden:eq:Li8}, and approximately $10^{-190}$ at $200$ working
digits for \eqref{golden:eq:Li9}; both coefficient vectors are primitive.
The weight-$9$ vector was recovered in three searches with different
padding arguments, all of which received coefficient $0$.
These observations support the formulas numerically. A residual printed
as zero at finite precision does not establish exact equality.
\begin{remark}
These experiments support golden formulas through weight $9$, using
four or five arguments per index at weight $5$, six arguments at weights
$6$--$7$, and eight arguments at weights $8$--$9$.
Higher-order golden extrapolations and their numerical construction
appear in Abouzahra--Lewin~\cite{golden:AbouzahraLewin1985} and
Bailey--Broadhurst~\cite{golden:BaileyBroadhurst1999}.
The present numerical evidence does not prove a maximal ladder weight.
Completeness of a multiplicative seed list, additional higher-weight
compatibility conditions, and numerical period independence are separate
questions.
\end{remark}

\section{Reducing the tetralogarithm relation to a functional identity}
Duplication reduces the proved tetralogarithm relation
\eqref{golden:eq:tetra} to the following equivalent cross-base relation.
The complete certificate proof that follows establishes both;
no independence claim about the six individual values is needed.

Two duplications
$\Li_4(\tfrac14)=8\bigl(\Li_4(\tfrac12)+\Li_4(-\tfrac12)\bigr)$ and
$\Li_4(\tfrac19)=8\bigl(\Li_4(\tfrac13)+\Li_4(-\tfrac13)\bigr)$ reduce the left side to the reflected arguments,
giving the following \emph{cross-base Kummer specialization}:
{\footnotesize
\begin{multline}\label{golden:eq:kummer}
  \Li_4(\tfrac34)+4\Li_4(\tfrac23)=-\tfrac{16}3\Li_4(\tfrac12)-4\Li_4(-\tfrac12)+2\Li_4(\tfrac13)-2\Li_4(-\tfrac13)
  +\tfrac{19}{540}\pi^4+\tfrac59\pi^2\ln^2 2+\tfrac16\pi^2\ln^2 3\\
  -\tfrac23\pi^2\ln2\ln3-\tfrac{19}{18}\ln^4 2-\tfrac16\ln^4 3+\tfrac43\ln^3 2\ln3-\ln^2 2\ln^2 3+\tfrac23\ln2\ln^3 3.
\end{multline}}
Substituting the two duplication identities into
\eqref{golden:eq:kummer} cancels the tetralogarithms in
\eqref{golden:eq:tetra}. The next proof supplies rational-function
cancellation, real branch conventions, and endpoint descent directly,
closing the earlier unproved-specialization gap.

% Proof notes for the root article. Requires amsmath, amssymb, amsthm,
% booktabs, hyperref. The table input is supplied beside the certificate.
% This is a complete mathematical contribution, not a proposed proof.

\section{An exact proof of the rational tetralogarithm conjecture}
\label{rattetra:sec}

The six-term relation \eqref{golden:eq:tetra} was proposed in 2015 and
2017~\cite{rattetra:2015,rattetra:2017}. Its numerical discovery and the
duplication reduction suggest a rational-function certificate. The
continuation~\cite{rattetra:report} supplies the forty-two rows below, together
with the analytic descent and endpoint calculations. These jointly prove the
relation; tensor cancellation alone would not fix the constants or branches.
The proof uses elementary differential identities for classical
polylogarithms and finite rational arithmetic, and also yields a
one-parameter functional identity. No worldwide priority or period
independence claim is needed.

\begin{theorem}[The six rational arguments]
\label{rattetra:thm:main}
Put $A=\log2$ and $B=\log3$. Then
\begin{align}
&720\operatorname{Li}_4(\tfrac12)
 -2160\operatorname{Li}_4(\tfrac13)
 +2160\operatorname{Li}_4(\tfrac23)
 +270\operatorname{Li}_4(\tfrac14)\notag\\
&\qquad+540\operatorname{Li}_4(\tfrac34)
 +135\operatorname{Li}_4(\tfrac19)\notag\\
&=19\pi^4+30(\pi^2-B^2)(10A^2-12AB+3B^2)
       -30A^2(19A^2-24AB+8B^2).
\label{rattetra:eq:main}
\end{align}
In particular, the intermediate relation \eqref{golden:eq:kummer}
above is also proved, by the two elementary duplications
already recorded there.
\end{theorem}

\subsection{A real generalized Rogers function and its exact descent}

Write $\ell_n(x)=\operatorname{Re}\operatorname{Li}_n(x)$ on the
principal branch. For $x\in\mathbb R\setminus\{0,1\}$ set
$L=\log|x|$, $M=\log|1-x|$, so $\ell_1(x)=-M$. Define
\begin{align*}
R_2(x)&=\ell_2(x)-\tfrac12L\ell_1(x),\\
R_3(x)&=\ell_3(x)-L\ell_2(x)+\tfrac13L^2\ell_1(x),\\
R_4(x)&=\ell_4(x)-L\ell_3(x)
                +\tfrac12L^2\ell_2(x)-\tfrac18L^3\ell_1(x).
\end{align*}
These functions extend continuously by $R_n(0)=0$ and
$R_n(1)=\zeta(n)$. Direct telescoping differentiation gives
\begin{equation}
 dR_n(x)=\frac{(-1)^n(n-1)}{n!}
 L^{n-2}(L\,dM-M\,dL),\qquad n=2,3,4.
\label{rattetra:eq:derivative}
\end{equation}
At a crossing of $x=1$, all logarithmically singular terms have positive
powers of $L$, which tend to zero; this proves the stated continuity and
allows identities obtained on regular subintervals to be joined.

Let $V=\mathbb Q(t)^\times\otimes_{\mathbb Z}\mathbb Q$, written
additively, and put
\[
 \beta_n[f]=f^{\otimes(n-2)}\otimes(f\wedge(1-f)).
\]
The sign of a rational function is torsion and vanishes in $V$.
Rational prime factors do not vanish and must be retained.

\begin{lemma}[Descent with several independent logarithms]
\label{rattetra:lem:descent}
Let $f_i(t)$ be nonconstant rational functions with no zero or pole
on an interval $I$, and let $v_i\in\mathbb Q$ satisfy
$\sum_i v_i\beta_4[f_i]=0$. For every linear functional
$\eta:V\to\mathbb R$ and each $n=2,3,4$,
\[
 \sum_i v_i\eta(f_i)^{4-n}R_n(f_i(t))
\]
is constant on $I$, with the continuous interpretation at $f_i(t)=1$.
\end{lemma}

\begin{proof}
Contract $4-n$ of the first two tensor slots with $\eta$.
The result is $\sum_i v_i\eta(f_i)^{4-n}\beta_n[f_i]=0$.
Evaluate the remaining ordinary tensor slots by $g\mapsto\log|g(t)|$
and the wedge by
$g\wedge h\mapsto\log|g|\,d\log|h|-\log|h|\,d\log|g|$.
Equation~\eqref{rattetra:eq:derivative} gives a zero derivative on every
regular subinterval. Continuity joins the constants across the finitely
many crossings of $f_i=1$.
\end{proof}

The point of the lemma is that $\eta(f_i)$ need not be a rational
multiple of a single logarithm. At the final specialization $t=1/2$,
the needed values are independent linear forms in $A=\log2$ and
$B=\log3$. Thus the one-logarithm golden descent does not have to be
applied outside its stated hypotheses.

\subsection{The complete finite certificate}

For the following forty-two rows define
\begin{equation}
 f_i(t)=\varepsilon_i2^{d_i}t^{a_i}(1-t)^{b_i}(1+t)^{c_i}.
\label{rattetra:eq:arguments}
\end{equation}
All $a_i\ge0$, and all factors except the displayed sign are positive
for $0<t<1$. The coefficients below are primitive integers.

\begin{center}\small
\setlength{\tabcolsep}{4pt}
\begin{tabular}{r r c r r r r@{\qquad}r r c r r r r}
\toprule
$i$&$v_i$&$\varepsilon_i$&$d_i$&$a_i$&$b_i$&$c_i$&
$i$&$v_i$&$\varepsilon_i$&$d_i$&$a_i$&$b_i$&$c_i$\\
\midrule
1 & -1 & + & -1 & 0 & 0 & 1 & 22 & 3 & + & -1 & 1 & 0 & -1 \\
2 & -11 & + & 0 & 0 & 0 & 1 & 23 & -4 & + & 0 & 1 & 0 & -1 \\
3 & -6 & - & 0 & 0 & 0 & 1 & 24 & 6 & - & 0 & 1 & 0 & -1 \\
4 & -3 & + & 1 & 0 & 0 & 1 & 25 & 1 & + & 1 & 1 & 0 & -1 \\
5 & -4 & + & 0 & 0 & 1 & -1 & 26 & -3 & + & -1 & 1 & 0 & 0 \\
6 & 4 & - & 0 & 0 & 1 & -1 & 27 & -3 & - & -1 & 1 & 0 & 0 \\
7 & -1 & + & -1 & 0 & 1 & 0 & 28 & 6 & + & 0 & 1 & 0 & 0 \\
8 & 22 & + & 0 & 0 & 1 & 0 & 29 & -6 & - & 0 & 1 & 0 & 0 \\
9 & 21 & - & 0 & 0 & 1 & 0 & 30 & 3 & + & 1 & 1 & 0 & 0 \\
10 & -3 & + & 1 & 0 & 1 & 0 & 31 & 3 & - & 1 & 1 & 0 & 0 \\
11 & 4 & + & 0 & 0 & 1 & 1 & 32 & 1 & + & -1 & 1 & 0 & 1 \\
12 & -1 & + & 0 & 0 & 2 & 1 & 33 & 1 & - & -1 & 1 & 1 & 0 \\
13 & -1 & - & 0 & 0 & 3 & 0 & 34 & 3 & + & 0 & 1 & 1 & 0 \\
14 & -1 & + & -1 & 1 & -2 & 0 & 35 & 3 & - & 0 & 1 & 1 & 0 \\
15 & -6 & - & 0 & 1 & -2 & 0 & 36 & -4 & - & 0 & 2 & -1 & -1 \\
16 & -6 & - & 0 & 1 & -1 & -1 & 37 & -3 & - & 0 & 2 & -1 & 0 \\
17 & 3 & - & -1 & 1 & -1 & 0 & 38 & -1 & + & 1 & 2 & -1 & 0 \\
18 & -21 & + & 0 & 1 & -1 & 0 & 39 & -3 & + & 0 & 2 & 0 & -1 \\
19 & -7 & - & 0 & 1 & -1 & 0 & 40 & -1 & + & 1 & 2 & 0 & -1 \\
20 & 1 & - & 1 & 1 & -1 & 0 & 41 & 1 & + & 0 & 3 & -3 & 0 \\
21 & -1 & - & -1 & 1 & 0 & -2 & 42 & 1 & - & 0 & 3 & -1 & -2 \\
\bottomrule
\end{tabular}
\end{center}

\begin{lemma}[Exact tensor cancellation]
\label{rattetra:lem:certificate}
The displayed rational functions satisfy
\begin{equation}
 \sum_{i=1}^{42}v_i\beta_4[f_i]=0.
\label{rattetra:eq:tensor}
\end{equation}
\end{lemma}

\begin{proof}
The complete factor list needed for all $f_i$ and $1-f_i$, modulo
sign, is
\[
 2,\quad t,\quad t-1,\quad t+1,\quad
 t-2,\quad t+2,\quad2t-1,\quad2t+1,\quad
 t^2-t-1,\quad t^2-t+1.
\]
Use their exponent vectors as coordinates in $V$. For example,
the constant $2$ occurs in
\[
 1-\frac{1+t}{2}=\frac{1-t}{2},\qquad
 1-\frac{t}{2(1-t)^2}
   =\frac{(2t-1)(t-2)}{2(1-t)^2}.
\]
The first two slots of $\beta_4$ are symmetric. For each of the ten
unordered pairs of coordinates of $2,t,1-t,1+t$, multiply the
corresponding exponents and the two-by-two alternating minors of
the vectors of $f_i,1-f_i$, and sum with coefficient $v_i$.
Every resulting coordinate is zero. There are $230$ coordinates
touched by the finite sum.

The accompanying \src{tetralogarithm_certificate.json} gives
exactly the forty-two displayed rows. The self-contained verifier
independently factors every rational function, checks the polynomial
factorizations by multiplication, retains the prime-$2$ coordinate,
and performs the indicated sums over rational integers. It does not
search for a relation or use numerical polylogarithms in this check.
This is a finite exact verification of~\eqref{rattetra:eq:tensor}.
\end{proof}

\subsection{The functional identity proved by the certificate}

For $0<t<1$, put
\[
 A=\log2,\qquad X=\log t,\qquad
 Y=\log(1-t),\qquad Z=\log(1+t),
\]
and define
\begin{align}
P_4={}&\frac{A^4}{4}-\frac{A^3(Y+Z)}6
       +\frac{A^2(Y^2+Z^2)}4+\frac{5A(Y^3+Z^3)}6
       -XY(3Y+Z)(Y+Z)\notag\\
 &+\frac{17Y^4+8Y^3Z+9Y^2Z^2+2YZ^3+2Z^4}{6},
 \label{rattetra:eq:P4}\\
P_2={}&-6A^2-4A(Y+Z)+\frac52Y^2+8YZ-8Z^2.
 \label{rattetra:eq:P2}
\end{align}

\begin{theorem}[A rational-function identity]
\label{rattetra:thm:functional}
With the forty-two rows of the table,
\begin{equation}
 \sum_{i=1}^{42}v_i\operatorname{Re}\operatorname{Li}_4(f_i(t))
 -4\operatorname{Li}_4(\tfrac12)
 =P_4+\zeta(2)P_2-\frac{71}{4}\zeta(4),
 \qquad 0<t<1.
\label{rattetra:eq:functional}
\end{equation}
The principal real parts and continuous values are used when an
argument is greater than one or equals one.
\end{theorem}

\begin{proof}
Fix $t\in(0,1)$ and write $L_i=\log|f_i(t)|$.
Apply Lemma~\ref{rattetra:lem:descent} with the functional sending a
prime or irreducible factor $2,s,1-s,1+s$ respectively to
$A,X,Y,Z$, and extend it linearly to the factor space, assigning
arbitrary values to its other basis elements. Unique factorization
shows that such a functional exists, and its value on every
$f_i(s)$ is $L_i$. This definition remains valid when one of the
additional factors of $1-f_i(t)$ vanishes; no logarithm of zero
is used to define the functional. The table functions have limits
\[
 e_i=\lim_{s\to0^+}f_i(s)=
 \begin{cases}0,&a_i>0,\\ \varepsilon_i2^{d_i},&a_i=0.
 \end{cases}
\]
Thus for $n=2,3,4$,
\[
 \sum_i v_iL_i^{4-n}R_n(f_i(t))
  =\sum_{a_i=0}v_iL_i^{4-n}R_n(e_i).
\]
The elementary triangular identity
\[
 \ell_4(x)=R_4(x)+L R_3(x)+\tfrac12L^2R_2(x)
                         +\tfrac1{24}L^3\ell_1(x)
\]
then supplies an exact ordinary-polylogarithm identity. More explicitly,
put $M_i=\log|e_i|=d_iA$ for $a_i=0$, and set
$Q_i=-M_i^3/8+L_iM_i^2/3-L_i^2M_i/4$. The result is
\begin{align}
\sum_i v_i\ell_4(f_i(t))={}&\sum_{a_i=0}v_i\bigl\{
 \ell_4(e_i)+(L_i-M_i)\ell_3(e_i)
 +\tfrac12(L_i-M_i)^2\ell_2(e_i)+Q_i\ell_1(e_i)\bigr\}\notag\\
 &+\frac1{24}\sum_i v_iL_i^3\ell_1(f_i(t)).
\label{rattetra:eq:descent-master}
\end{align}
When $e_i=1$, $Q_i=0$ and the corresponding last term in braces
is zero; no divergent value is assigned to $\ell_1(1)$.
Likewise $L_i^3\ell_1(f_i(t))$ is assigned its limiting value zero
at $f_i(t)=1$.

Only $e_i\in\{1,-1,1/2,2\}$ occur. The required lower values are
\[
\begin{array}{c|ccc}
e&\ell_1(e)&\ell_2(e)&\ell_3(e)\\ \hline
-1&-A&-\zeta(2)/2&-3\zeta(3)/4\\
1/2&A&\zeta(2)/2-A^2/2&7\zeta(3)/8-\zeta(2)A/2+A^3/6\\
2&0&3\zeta(2)/2&7\zeta(3)/8+3\zeta(2)A/2
\end{array}
\]
and $\ell_2(1)=\zeta(2)$, $\ell_3(1)=\zeta(3)$.
For completeness, the lower reflection and Landen identities are
\begin{align*}
\operatorname{Li}_2(x)+\operatorname{Li}_2(1-x)
 &=\zeta(2)-\log x\log(1-x),\\
\operatorname{Li}_2\!\left(-\frac{x}{1-x}\right)
 &=-\operatorname{Li}_2(x)-\frac12\log^2(1-x),\\
\operatorname{Li}_3(x)+\operatorname{Li}_3(1-x)
 +\operatorname{Li}_3\!\left(-\frac{x}{1-x}\right)
 &=\zeta(3)+\zeta(2)\log(1-x)\\
 &\quad-\frac12\log x\log^2(1-x)+\frac16\log^3(1-x),
\end{align*}
for $0<x<1$. Differentiating the first two identities gives zero
differences, with constants fixed at $x\to0^+$. Using them in the
derivative of the third gives zero as well, and the same endpoint
fixes its constant. Substitution $x=1/2$ proves the displayed half
values. The values at $2$ follow by real inversion. Thus these
endpoint reductions require no integer-relation recognition.

The endpoint weight-four sum is exactly
\[
 \sum_{a_i=0}v_i\ell_4(e_i)
 =4\operatorname{Li}_4(\tfrac12)+\frac{A^4}{4}
       -6\zeta(2)A^2-\frac{71}{4}\zeta(4).
\]
Substitute $L_i=d_iA+a_iX+b_iY+c_iZ$ and the factorizations of
$1-f_i$ in~\eqref{rattetra:eq:descent-master}. All logarithms of the
six additional polynomial factors cancel, as do the $\zeta(3)$ terms.
The remaining polynomial is precisely
$P_4+\zeta(2)P_2-71\zeta(4)/4$, which proves the theorem.
The verifier expands both the master formula and the displayed
$P_4,P_2$ independently and checks their exact polynomial difference
is zero.
\end{proof}

\subsection{The rational specialization, including all inversion terms}

For $r>1$, the real inversion formulas at weight four are
\begin{align*}
 \ell_4(r)+\operatorname{Li}_4(r^{-1})
   &=-\frac{\log^4r}{24}+\zeta(2)\log^2r+2\zeta(4),\\
 \operatorname{Li}_4(-r)+\operatorname{Li}_4(-r^{-1})
   &=-\frac{\log^4r}{24}-\frac12\zeta(2)\log^2r
       -\frac74\zeta(4).
\end{align*}
They follow by repeated differentiation from the real logarithm
and dilogarithm inversion laws, fixing the constants at $r=1$.
Negative arguments inside the unit interval are removed by
\[
 \operatorname{Li}_4(-r)=\tfrac18\operatorname{Li}_4(r^2)
                                  -\operatorname{Li}_4(r),
 \qquad 0\le r\le1.
\]
At $t=1/2$ one has $X=Y=-A$, $Z=B-A$ and
$f_i(1/2)=\varepsilon_i2^{d_i-a_i-b_i-c_i}3^{c_i}$.
Applying the three displayed rules to the left side of
\eqref{rattetra:eq:descent-master}, after subtracting its endpoint
weight-four sum and multiplying by $180$, gives exactly the six
tetralogarithms in~\eqref{rattetra:eq:main}. The remaining polynomial
from those inversions is
\begin{align*}
I={}&120A^4-510A^3B+765A^2B^2-510AB^3+150B^4\notag\\
 &+\zeta(2)(90A^2+2880AB-1980B^2)-1710\zeta(4).
\end{align*}
The non-weight-four part of the specialized descent master formula is
\begin{align*}
D={}&-450A^4+210A^3B+225A^2B^2-150AB^3+60B^4\notag\\
 &+\zeta(2)(1890A^2+720AB-1440B^2).
\end{align*}
Consequently the six-term sum equals $D-I$, namely
\begin{align*}
&-570A^4+720A^3B-540A^2B^2+360AB^3-90B^4\notag\\
&\qquad+\zeta(2)(1800A^2-2160AB+540B^2)+1710\zeta(4).
\end{align*}
Using $\zeta(2)=\pi^2/6$, $\zeta(4)=\pi^4/90$ gives exactly
\eqref{rattetra:eq:main}, proving Theorem~\ref{rattetra:thm:main}.

% Optional insertion after identities_notes.tex; the certificate is shared.
% Requires the same theorem environments, amsmath and booktabs.
\subsection{A further proved rational specialization}
\label{rattetra:sec:onethird}

The functional identity also produces explicit identities beyond the
six-term conjecture. The specialization $t=1/3$ is convenient because
all logarithms reduce to $A=\log2$ and $B=\log3$. The following
corollary is a new entry in this report, with no claim of priority
among classical tetralogarithm identities.

\begin{corollary}[Twenty-five positive rational arguments]
\label{rattetra:cor:onethird}
For the argument--coefficient pairs $(r,c_r)$ in the table,
\begin{align}
\sum_r c_r\operatorname{Li}_4(r)
={}&\frac{109}{3}A^4-116A^3B+54A^2B^2+36AB^3-21B^4\notag\\
 &+\zeta(2)(316A^2-408AB+132B^2)+53\zeta(4).
\label{rattetra:eq:onethird}
\end{align}
All arguments lie in $(0,1)$, and the integer coefficient list is
primitive: its greatest common divisor is one.

\begin{center}
\renewcommand{\arraystretch}{1.25}
\begin{tabular}{c r@{\qquad\qquad}c r}
\toprule
$r$&$c_r$&$r$&$c_r$\\
\midrule
$1/1024$ & $1$ & $9/64$ & $-6$ \\
$9/1024$ & $-1$ & $1/6$ & $16$ \\
$1/81$ & $1$ & $2/9$ & $8$ \\
$1/64$ & $-4$ & $1/4$ & $-107$ \\
$1/36$ & $-6$ & $8/27$ & $8$ \\
$1/32$ & $-8$ & $1/3$ & $80$ \\
$4/81$ & $3$ & $3/8$ & $64$ \\
$1/16$ & $9$ & $4/9$ & $24$ \\
$1/12$ & $-24$ & $1/2$ & $-200$ \\
$64/729$ & $-1$ & $16/27$ & $-8$ \\
$3/32$ & $8$ & $3/4$ & $112$ \\
$1/9$ & $-14$ & $8/9$ & $32$ \\
$1/8$ & $64$ & & \\
\bottomrule
\end{tabular}
\end{center}
\end{corollary}

\begin{proof}
Set $t=1/3$ in Theorem~\ref{rattetra:thm:functional}. The logarithms
are $X=-B$, $Y=A-B$, and $Z=2A-B$. Apply the real inversion
formulas to arguments whose absolute value exceeds one, then the
duplication formula
$\operatorname{Li}_4(-r)=\operatorname{Li}_4(r^2)/8-
\operatorname{Li}_4(r)$ to negative arguments in $[-1,0]$.
Multiplication by eight clears the resulting coefficient denominators.
Collecting equal positive arguments gives exactly the table, and
collecting all inversion terms gives the displayed polynomial.
The same certificate verifier performs these reductions over
$\mathbb Q$ and checks that the ordinary-polynomial residual is zero.
As an independent audit only, a direct evaluation of the reduced
identity at ninety working decimal digits has absolute residual
below $8\times10^{-90}$.
\end{proof}

The argument list and polynomial are supplied separately in
\src{tetralogarithm_one_third.json}, and are also included in the
shared certificate data. The replay checks the original conjecture,
the functional identity, and this corollary in a single run.

\subsection{Verification, provenance, and further questions}

The exact verifier finishes with a zero tensor residual, a zero
ordinary-polynomial residual, and the six specified rational arguments
with coefficients $(720,-2160,2160,270,540,135)$. Separately, it evaluates
the functional identity at $t=1/5,37/100,1/2,4/5$ using $90$ working
decimal digits. The largest numerical discrepancy is approximately
$7.9\times10^{-90}$. These calculations check the branch conventions;
the proof is the finite tensor cancellation and the analytic descent.

The original questions are
\url{https://math.stackexchange.com/questions/1373123/} and
\url{https://mathoverflow.net/questions/261408/}.
The rational weight-four question and the separate golden weight-five
ladders have different proof obligations: the certificate here settles the rational
weight-four formula directly. The method is compatible with the
higher-Bloch-condition approach to functional equations, but no unproved
period conjecture or transcendence assertion is used.

Three useful next problems are to decompose the forty-two rows into a
small set of named classical Kummer specializations; to specialize the
proved functional identity at other rational or quadratic arguments and
systematically reduce the resulting baskets; and to search for similarly
complete certificates for the unresolved golden weights six and seven.
The Gaussian $S_6$ conjecture remains a distinct problem and is not
promoted by this result. No minimality or arithmetic independence of
the six rational tetralogarithms is asserted.

The complete delivered certificate and verifier are preserved under
\src{reports/uniform-transition-continuation/continuations/tetralogarithm-euler-polynomial-zeros/code/}.
The fresh isolated replay is recorded under
\src{verification/seventh-replay/tetra/}. It checks all forty-two rows and
230 touched tensor coordinates, the exact lower-weight polynomial descent,
both rational specializations, and separately labeled numerical branch
diagnostics. The other claims in that package have independent proofs and
retain their editorial-ledger audit status.


\section{Canonical index-12/20/24 ladders and a $\zeta$-elimination chain}
Reshetnikov's $\Li_5$ conjectures are the weight-$5$ entries of three first-degree golden ladders (Abouzahra--Lewin).
For positive integer $n$, use the following canonical
(monic, $a^{n-1}$-normalized) notation. Put $\rho=\phi^{-1}$,
$\ell=\ln\rho$, and
\[
\mathrm{tw}(n)=D_1\frac{\ell^n}{n!}
 +D_2\zeta(2)\frac{\ell^{n-2}}{(n-2)!}
 +D_3\zeta(4)\frac{\ell^{n-4}}{(n-4)!},
\]
where a summand with a negative factorial index is omitted.
The previously implicit $L_{12}$ is reconstructed from the first
displayed weight-$5$ formula using the same normalization:
\begin{align*}
L_{12}(n)&=\frac{\Li_n(\rho^{12})}{12^{n-1}}
-\frac32\frac{\Li_n(\rho^6)}{6^{n-1}}
-\frac{\Li_n(\rho^4)}{4^{n-1}}
+\frac{11}{48}\frac{\Li_n(\rho^2)}{2^{n-1}}+\mathrm{tw}(n),\\
L_{20}(n)&=\frac{\Li_n(\rho^{20})}{20^{n-1}}
-\frac{\Li_n(\rho^{10})}{10^{n-1}}-3\frac{\Li_n(\rho^4)}{4^{n-1}}
+\frac12\frac{\Li_n(\rho^2)}{2^{n-1}}+\mathrm{tw}(n),\\
L_{24}(n)&=\frac{\Li_n(\rho^{24})}{24^{n-1}}
-\frac43\frac{\Li_n(\rho^{12})}{12^{n-1}}
-2\frac{\Li_n(\rho^8)}{8^{n-1}}\\
&\quad+\frac73\frac{\Li_n(\rho^4)}{4^{n-1}}
-\frac{205}{576}\frac{\Li_n(\rho^2)}{2^{n-1}}+\mathrm{tw}(n).
\end{align*}
For $L_{12},L_{20},L_{24}$ respectively, the triples $(D_1,D_2,D_3)$ are
\[
(-13/48,1/48,-19/1728),\quad
(-1/2,1/25,-89/4000),\quad
(179/576,-5/192,629/41472).
\]
The displayed weight-$5$ formulas are equivalent, by exact rational
recombination and normalization, to
\[
\begin{aligned}
L_{12}(5)&=\frac{67}{6912}\zeta(5),&
L_{20}(5)&=\frac{201}{10000}\zeta(5),\\
L_{24}(5)&=-\frac{1541}{110592}\zeta(5).
\end{aligned}
\]
For the last formula, first add $64/3$ times the first displayed
weight-$5$ formula to the third, eliminating $\Li_5(\rho^6)$, then
divide by $-15\cdot24^4$.
This supplies the missing notation; it does not provide a new proof of
those evaluations. The source's low-order vanishing statements, with the
factorial convention above, concern $1\le n\le4$.
The rational combinations
$M_{20}=L_{20}-\tfrac{1296}{625}L_{12}$ and
$M_{24}=L_{24}+\tfrac{23}{16}L_{12}$ cancel the displayed
$\zeta(5)$ coefficients exactly. The source's subsequent numerical
extrapolation reaches weight $7$ and forms
$T=M_{24}+\tfrac{409375}{373248}M_{20}$ to cancel its
$\zeta(7)$ term and reach weight $9$, proposing
\[
T(9)=\frac{66452911}{41278242816000}\zeta(9)
 +\frac{3537283}{2063912140800}\zeta(8)\ell
 -\frac{11527}{2866544640}\zeta(6)\frac{\ell^3}{3!}.
\]
The source reports direct substitution and an independent PSLQ check at
$460$ digits for this expression. Those checks record numerical evidence
for the higher-order extrapolation, rather than a functional-equation
certificate. This presentation follows the Abouzahra--Lewin ladder
construction~\cite{golden:AbouzahraLewin1985}; the single eight-argument
formulas \eqref{golden:eq:Li8}--\eqref{golden:eq:Li9} are written separately.

\section{The minimal Pisot ladder base}
Beyond the golden ratio, the reciprocal of the minimal Pisot number
$\omega=0.7548776\ldots$, the real root of $x^3+x^2-1$, carries its
own ladders. Abouzahra, Lewin, and Xiao Hongnian studied this field in
their published 1987 sequel~\cite{golden:AbouzahraLewinXiao1987}.
Its seeds include the pure-unit relations
$1-\omega=\omega^5$, $1-\omega^2=\omega^3$,
$1-\omega^3=\omega^2$, and $1-\omega^5=\omega$.
The source records the following low-weight formulas, with
$\lambda=\ln\omega$:
{\footnotesize
\begin{align*}
  &\tfrac12\Li_2(\omega^2)-\Li_2(\omega^5)-\lambda^2=0,\qquad
  -\Li_3(\omega)+\tfrac64\Li_3(\omega^2)+\tfrac39\Li_3(\omega^3)-\Li_3(\omega^5)-\tfrac{6\lambda^3}{3!}=0,\\
  &-71\Li_4(\omega)-\tfrac{120}8\Li_4(\omega^2)+\tfrac{1053}{27}\Li_4(\omega^3)-17\Li_4(\omega^5)+\tfrac{5488}{343}\Li_4(\omega^7)-\Li_4(\omega^{14})+51\zeta(4)+\tfrac{2010\lambda^4}{4!}-\tfrac{84\zeta(2)\lambda^2}{2!}=0.
\end{align*}}
The weight-four formula uses the extra legs $\omega^7,\omega^{14}$.
The later Bloch-framework investigation identifies higher ladders for this
field; the earlier assertion that the tower stops at weight four is
discarded. A fixed argument set can stop yielding detected relations while
larger sets continue to produce them.

\section{Hunting more ladder bases, and a methodological caveat}
Pure-unit seeds $1-\rho^a=\rho^b$ offer a useful way to discover ladders.
They do not exhaust general multiplicative relations, and a bounded search
over polynomial coefficients does not cover all monic polynomials of a
given degree. Root disguises are easy to miss: the root of $x^4+x^2-1$
becomes a golden reciprocal after squaring. The following two cubic and
quartic relations illustrate why negative searches require this care.

Two bases the earlier survey had dismissed as ``weight-$2$ only'' do, in fact, carry weight-$3$ ladders. With
$\rho$ the root of $x^3+x-1$ (the supergolden reciprocal) and of $x^4+x-1$ (the reciprocal of the second-smallest
Pisot number $x^4-x^3-1$),
{\footnotesize
\begin{align*}
  -60\Li_3(\rho)+54\Li_3(\rho^2)-32\Li_3(\rho^3)-3\Li_3(\rho^6)&=24\ln^3\rho-4\pi^2\ln\rho-36\zeta(3),\\
  -12\Li_3(\rho)+12\Li_3(\rho^3)-12\Li_3(\rho^4)-3\Li_3(\rho^6)&=22\ln^3\rho-2\pi^2\ln\rho-12\zeta(3),
\end{align*}}
\!each verified in two engines (and by direct substitution to $400$+ digits). The lesson is methodological:
\textbf{a PSLQ ``no relation'' is not a proof of non-existence.} The first relation above was \emph{missed} by a
pure-tower PSLQ search at $300$ digits, yet passed direct substitution
at $400$ digits. Its exact proof below establishes that this was a false
negative. Recomputing a frozen candidate at precision well above its
discovery precision is a useful diagnostic; the analytic proofs below
establish the two identities. Whether these two bases climb past weight
$3$ remains open in the present investigation; the question is delicate
precisely because PSLQ negatives cannot settle it.
% Self-contained contribution for the Polylogarithms continuation article.
% Required packages: amsmath, amssymb, amsthm, booktabs, hyperref.
% Required theorem environments: theorem, lemma, corollary, remark.

\section{Proof certificates for the cubic and quartic trilogarithm ladders}
\label{gaussian:sec:proved-small-ladders}

The two displayed formulas record two trilogarithm combinations
at the positive roots of $r^3+r=1$ and $r^4+r=1$, with substantial experimental evidence. Their exact functional-equation
certificates now explain the coefficients.
We give such derivations. The quartic identity belongs to a uniform
family; the cubic identity has a short certificate using two
specializations of Kummer's nine-term relation and four Rogers pentagons.
These are exact consequences of classical functional equations, not
claims that the underlying equations are new.

\subsection{A uniform ladder from a complementary pair of powers}

For $0<x<1$ the real Landen relation for the trilogarithm is
\begin{align}
\operatorname{Li}_3(x)+\operatorname{Li}_3(1-x)
 +\operatorname{Li}_3(1-x^{-1})
 &=\zeta(3)+\frac{\pi^2}{6}\log x
 -\frac12\log^2x\log(1-x)+\frac16\log^3x.
\label{gaussian:eq:real-trilog-landen}
\end{align}
For a modern account and its arithmetic analogues see
Nakamura and Shiraishi, \emph{Landen's trilogarithm functional equation
and $\ell$-adic Galois multiple polylogarithms} \cite{gaussian:NakamuraShiraishi}.
For completeness, a real-variable proof is especially short.
Differentiation of the left side, followed by
\[
\operatorname{Li}_2(x)+\operatorname{Li}_2(1-x)
 =\frac{\pi^2}{6}-\log x\log(1-x),\qquad
\operatorname{Li}_2(1-x^{-1})
 =-\operatorname{Li}_2(1-x)-\frac12\log^2x,
\]
gives
\[
\frac{\pi^2/6-\log x\log(1-x)}{x}
  +\frac{\log^2x}{2x(1-x)}.
\]
This is the derivative of the right side. The difference tends to zero
as $x\to1^-$, proving \eqref{gaussian:eq:real-trilog-landen}. The two dilogarithm
relations used here follow themselves by one differentiation and an
endpoint value. All logarithms in the calculation are real, so there
is no branch ambiguity.

\begin{theorem}[Complementary-power trilogarithm ladder]
\label{gaussian:thm:complementary-powers}
Let $a,b$ be positive integers with $a\le b$, and let $0<\rho<1$ be
the unique positive solution of $\rho^a+\rho^b=1$. Set
$\lambda=\log\rho$ and $d=b-a$. Then
\begin{align}
&\operatorname{Li}_3(\rho^a)+\operatorname{Li}_3(\rho^b)
 -\operatorname{Li}_3(\rho^d)
 +\frac14\operatorname{Li}_3(\rho^{2d})\notag\\
&\hspace{20mm}=\zeta(3)+\frac{a\pi^2}{6}\lambda
 +\frac{a^2(a-3b)}{6}\lambda^3.
\label{gaussian:eq:complementary-powers}
\end{align}
When $a=b$, the powers $\rho^0$ are interpreted as $1$.
\end{theorem}

\begin{proof}
The polynomial $t^a+t^b$ is strictly increasing from $0$ to $2$ on
$[0,1]$, which proves existence and uniqueness. Substitute $x=\rho^a$
in \eqref{gaussian:eq:real-trilog-landen}; its other arguments are
$1-x=\rho^b$ and $1-x^{-1}=-\rho^{b-a}$. The absolutely convergent
series identity
\[
\operatorname{Li}_3(z)+\operatorname{Li}_3(-z)
 =\frac14\operatorname{Li}_3(z^2),\qquad 0\le z\le1,
\]
eliminates the negative argument. The logarithmic terms become exactly
the right side of \eqref{gaussian:eq:complementary-powers}.
\end{proof}

In particular, if $r_m^m+r_m=1$ and $\lambda_m=\log r_m$, then
\begin{align}
\operatorname{Li}_3(r_m)+\operatorname{Li}_3(r_m^m)
 -\operatorname{Li}_3(r_m^{m-1})
 +\frac14\operatorname{Li}_3(r_m^{2m-2})
 =\zeta(3)+\frac{\pi^2}{6}\lambda_m
 -\frac{3m-1}{6}\lambda_m^3.
\label{gaussian:eq:all-m-trilog}
\end{align}
This supplies an infinite family of finite exact reductions, although
different pairs $(a,b)$ can describe the same base after a power
substitution. It does not assert that these exhaust all ladders for
any one number field.

\begin{corollary}[The quartic identity]
If $r^4+r=1$ and $\lambda=\log r$, then
\begin{align}
-12\operatorname{Li}_3(r)+12\operatorname{Li}_3(r^3)
 -12\operatorname{Li}_3(r^4)-3\operatorname{Li}_3(r^6)
 =22\lambda^3-2\pi^2\lambda-12\zeta(3).
\label{gaussian:eq:quartic-proved}
\end{align}
\end{corollary}
\begin{proof}
Take $(a,b)=(1,4)$ in Theorem~\ref{gaussian:thm:complementary-powers}
and multiply by $-12$.
\end{proof}

For $m=2$, \eqref{gaussian:eq:all-m-trilog} reduces to the familiar golden value
\[
\operatorname{Li}_3(r_2^2)
 =\frac45\zeta(3)+\frac{2\pi^2}{15}\log r_2
 -\frac23\log^3r_2.
\]
For $m=3$ it gives a useful additional relation on the arguments
$r,r^2,r^3,r^4$; the manuscript's cubic identity instead uses
$r,r^2,r^3,r^6$. The change of the last argument requires more than
the one-variable Landen relation.

\subsection{An elementary dilogarithm certificate in the cubic field}

Throughout this subsection let $r^3+r=1$, $0<r<1$, and put
$\lambda=\log r$, $A=\log(1+r)$. Direct algebra in
$\mathbb Q[r]/(r^3+r-1)$ gives
\begin{align}
1-r&=r^3,&1-r^2&=r^3(1+r),&1-r^3&=r,\\
1-r^4&=r^2(1+r),&1+r^3&=r^2(1+r)^2,
 &1-r^6&=r^3(1+r)^2.
\label{gaussian:eq:cubic-factorizations}
\end{align}
Define the Rogers dilogarithm on $(0,1)$ by
\[
R(x)=\operatorname{Li}_2(x)+\frac12\log x\log(1-x).
\]
It satisfies $R(x)+R(1-x)=\zeta(2)$ and the pentagon
\begin{align}
R(x)+R(y)
 =R(xy)+R\!\left(\frac{x(1-y)}{1-xy}\right)
  +R\!\left(\frac{y(1-x)}{1-xy}\right),\qquad 0<x,y<1.
\label{gaussian:eq:rogers-pentagon}
\end{align}
The latter is a real-domain form of the classical five-term relation.
It may also be verified by differentiating in $x$ with fixed $y$ and
using the limit at $x=0$.

\begin{lemma}[Cubic dilogarithm identity]
\label{gaussian:lem:cubic-dilog}
\begin{equation}
\operatorname{Li}_2(r^6)
 =6\operatorname{Li}_2(r^2)-2\operatorname{Li}_2(r)
  -4\operatorname{Li}_2(r^3).
\label{gaussian:eq:cubic-dilog}
\end{equation}
\end{lemma}

\begin{proof}
Set $s=r/(1+r)$ and $t=r^4/(1+r)$, so every argument below lies
in $(0,1)$. Write $R_k=R(r^k)$ and $R_s=R(s)$, $R_t=R(t)$.
Applying \eqref{gaussian:eq:rogers-pentagon} at $(x,y)=(r,r)$ gives
\[
2R_1=R_2+2R_s.
\]
At $(x,y)=(r,r^3)$ its two remaining arguments are $1-s$ and $t$.
Since $R_1+R_3=\zeta(2)$, reflection gives
\[
R_t=R_s-R_4.
\]
At $(x,y)=(r^2,r^4)$ its two remaining arguments are $s$ and $t$,
and therefore
\[
R_6=R_2+2R_4-2R_s.
\]
Finally, at $(x,y)=(r,r^2)$ these arguments are $1-r^2$ and $r^4$.
Using reflection twice gives
\[
R_4=2R_2-2R_3.
\]
Eliminating $R_s,R_t,R_4$ now yields
$R_6=6R_2-2R_1-4R_3$. The difference between $R(r^k)$ and
$\operatorname{Li}_2(r^k)$ is
$(k\lambda/2)\log(1-r^k)$. In the indicated combination these
terms cancel: by \eqref{gaussian:eq:cubic-factorizations}, the required check is
\[
3\lambda(3\lambda+2A)
 -6\lambda(3\lambda+A)+3\lambda^2+6\lambda^2=0.
\]
This proves the claimed ordinary dilogarithm identity.
\end{proof}

\subsection{Two Kummer substitutions prove the cubic trilogarithm}

The single-valued trilogarithm used in the manuscript is
\[
P_3(z)=\operatorname{Re}\operatorname{Li}_3(z)
 -\log|z|\operatorname{Re}\operatorname{Li}_2(z)
 -\frac13\log^2|z|\log|1-z|.
\]
The value at $z=1$ is defined by continuity. It obeys
\[
P_3(z^{-1})=P_3(z),\qquad
P_3(z)+P_3(-z)=\frac14P_3(z^2),
\]
and the single-valued Landen relation
$P_3(x)+P_3(1-x)+P_3(1-x^{-1})=\zeta(3)$ for real $0<x<1$.
The latter follows from \eqref{gaussian:eq:real-trilog-landen} and the two
dilogarithm identities in its proof.

We use Kummer's classical equation in the form stated and justified
by Zagier, \emph{Special values and functional equations of
polylogarithms}, \S7, Example~3, equation~(7) \cite{gaussian:Zagier}:
\begin{align}
&2P_3(x)+2P_3(y)
 +2P_3\!\left(\frac{x(1-y)}{x-1}\right)
 +2P_3\!\left(\frac{y(1-x)}{y-1}\right)\notag\\
&\quad+2P_3\!\left(\frac{1-x}{1-y}\right)
 +2P_3\!\left(\frac{x(1-y)}{y(1-x)}\right)
 -P_3(xy)-P_3(x/y)\notag\\
&\quad-P_3\!\left(\frac{x(1-y)^2}{y(1-x)^2}\right)=2\zeta(3).
\label{gaussian:eq:kummer-P3}
\end{align}
We only use specializations with nonzero denominators and finite
arguments. Since $P_3$ is single-valued, some real arguments exceeding
$1$ cause no branch correction in this equation; after inversion and
duplication, all remaining arguments are $r^k\in(0,1)$.

\begin{theorem}[The cubic identity]
\label{gaussian:thm:cubic-proved}
Let $r^3+r=1$ with $0<r<1$, and let $\lambda=\log r$. Then
\begin{align}
&-60\operatorname{Li}_3(r)+54\operatorname{Li}_3(r^2)
 -32\operatorname{Li}_3(r^3)-3\operatorname{Li}_3(r^6)\notag\\
&\hspace{20mm}=24\lambda^3-4\pi^2\lambda-36\zeta(3).
\label{gaussian:eq:cubic-proved}
\end{align}
\end{theorem}

\begin{proof}
Write $P_k=P_3(r^k)$. Landen at $r$ and duplication give
\[
\mathsf A:=P_1-P_2+P_3+\tfrac14P_4=\zeta(3).
\]
The nine arguments of \eqref{gaussian:eq:kummer-P3} at
$(x,y)=(r,r^3)$ and $(-r,r^2)$, in their displayed order, are
\[
\begin{array}{c|rrrrrrrrr}
 &1&2&3&4&5&6&7&8&9\\ \hline
(r,r^3)&r&r^3&-r^{-1}&-r^5&r^2&r^{-4}&r^4&r^{-2}&r^{-6}\\
(-r,r^2)&-r&r^2&r^4&-r^{-1}&r^{-3}&-r^2&-r^3&-r^{-1}&-r^5
\end{array}
\]
respectively. Using inversion and
$P_3(-r^k)=\tfrac14P_{2|k|}-P_{|k|}$, these two equations reduce to
\begin{align*}
\mathsf B&:=\tfrac32P_2+2P_3+P_4-2P_5-P_6+\tfrac12P_{10}
 =2\zeta(3),\\
\mathsf C&:=-3P_1+\tfrac34P_2+3P_3+\tfrac52P_4+P_5
 -\tfrac14P_6-\tfrac14P_{10}=2\zeta(3).
\end{align*}
The exact rational row operation $4\mathsf C+2\mathsf B-48\mathsf A$
therefore proves
\begin{equation}
-60P_1+54P_2-32P_3-3P_6=-36\zeta(3).
\label{gaussian:eq:cubic-P3-certificate}
\end{equation}

To recover ordinary trilogarithms, put
$c_1=-60,c_2=54,c_3=-32,c_6=-3$ and
\[
D=\sum_{k\in\{1,2,3,6\}}c_k\operatorname{Li}_3(r^k),\quad
S=\sum_{k\in\{1,2,3,6\}}c_k k\operatorname{Li}_2(r^k),\quad
T=\sum_{k\in\{1,2,3,6\}}c_k k^2\log(1-r^k).
\]
Expansion of the definition gives
$\sum c_kP_k=D-\lambda S-\lambda^2 T/3$.
Lemma~\ref{gaussian:lem:cubic-dilog}, together with
$\operatorname{Li}_2(r)+\operatorname{Li}_2(r^3)
=\pi^2/6-3\lambda^2$, gives
\[
S=72\lambda^2-4\pi^2.
\]
The factorizations \eqref{gaussian:eq:cubic-factorizations} give
\[
T=-60(3\lambda)+216(3\lambda+A)
 -288\lambda-108(3\lambda+2A)=-144\lambda.
\]
Consequently \eqref{gaussian:eq:cubic-P3-certificate} becomes
$D-24\lambda^3+4\pi^2\lambda=-36\zeta(3)$, which is
\eqref{gaussian:eq:cubic-proved}.
\end{proof}

\subsection{What the certificate establishes and what remains open}

The proofs settle the two displayed cubic and quartic trilogarithm
identities without period-independence assumptions, numerical recognition,
or unproved higher Bloch-group assertions. The universal family
\eqref{gaussian:eq:complementary-powers} also gives a principled generator of
additional exact identities. It does not justify pseudointegration to
weight four or above: differentiating with respect to an algebraic base
does not preserve its isolated defining relation.

Natural next problems are to derive comparable local certificates for the
weight-five through weight-nine golden formulas, find the smallest useful
functional-equation closure for a chosen base, and extend the exact
argument/substitution checker to arbitrary algebraic number fields.
Any claim of minimality of the surviving period basket remains distinct
from the existence of the displayed reductions.

% Contribution for the continuing Polylogarithms report.
% Requires amsmath, amssymb, amsthm, booktabs, hyperref.
% Citation keys: Ljunggren1960, Tverberg1960, BoydMartinThom2015,
% AartsFokkinkKruijtzer2001.

\section{Rigidity of complementary powers}
\label{rigid:sec:complementary-rigidity}

The earlier complementary-power ladder starts from a relation
$r^a+r^b=1$, where $0<r<1$ and $a,b$ are positive integers.
That construction raises a concrete completeness question: can one base
have many different pairs of complementary powers?  We settle this
question for pure powers, including arbitrary common divisors in the
exponents.  The decisive arithmetic input is a classical irreducibility
theorem for trinomials.  Our conclusions concern this precisely specified
class of seeds; they do not classify general multiplicative relations
among $1-r^k$, or bound the maximum weight of polylogarithm ladders.

For $0<r<1$, define
\[
\mathcal C(r)=\{(a,b)\in\mathbb Z_{>0}^2:
                 a\le b,\ r^a+r^b=1\}.
\]
Let $\omega$ be the unique root in $(0,1)$ of
$X^3+X^2-1$.  Thus $\omega$ is the reciprocal of the plastic number.

\begin{theorem}[Complete classification of collisions]
\label{rigid:thm:complement-collisions}
For every real $0<r<1$, the set $\mathcal C(r)$ has at most two elements.
It has two elements if and only if $r^d=\omega$ for some positive
integer $d$, in which case
\[
\mathcal C(r)=\{(d,5d),(2d,3d)\}.
\]
If $\mathcal C(r)$ contains $(a,a)$, then
$\mathcal C(r)=\{(a,a)\}$ and $r=2^{-1/a}$.
\end{theorem}

The theorem supplies an exact answer to the seed-classification question
posed after the earlier cubic and quartic ladder certificates.  It is
proved below as a consequence of the Ljunggren--Tverberg theorem, with
all additional arguments included.  We make no claim that this
consequence is globally new in the literature.  The related theorem of
Aarts, Fokkink, and Kruijtzer classifies the two morphic numbers, by
requiring both $p+1$ and $p-1$ to be powers of $p$; their proof also uses
trinomial irreducibility \cite{rigid:AartsFokkinkKruijtzer2001}.  Our formulation
permits arbitrary positive exponent pairs and explicitly handles power
substitutions.

\subsection{The irreducibility input and the entire cyclotomic factor}

For $1\le a<b$, put $d=\gcd(a,b)$, $A=a/d$, $B=b/d$, and
\[
F_{a,b}(X)=X^b+X^a-1.
\]
The following precise specialization is the arithmetic input.

\begin{theorem}[Ljunggren--Tverberg, specialized]
\label{rigid:thm:trinomial-input}
The polynomial $F_{a,b}$ has exactly one irreducible factor over
$\mathbb Q$ whose roots are not roots of unity.  Its other factor is
\[
C_{a,b}(X)=
\begin{cases}
 X^{2d}-X^d+1,&(A,B)\equiv(1,5)\text{ or }(5,1)\pmod6,\\
 1,&\text{otherwise}.
\end{cases}
\]
In particular $H_{a,b}=F_{a,b}/C_{a,b}$ is irreducible, monic,
and has constant term $-1$.  Moreover
$H_{a,b}(X)=H_{A,B}(X^d)$.
\end{theorem}

This is the trinomial result of Ljunggren and Tverberg
\cite{rigid:Ljunggren1960,rigid:Tverberg1960}; a modern explicit formulation,
including the assertion that $H_{A,B}(X^d)$ remains irreducible, is
Boyd--Martin--Thom, Lemma~4.1(a),(b) \cite{rigid:BoydMartinThom2015}.
Only the trinomial theorem is used.  Its separate quadrinomial analogue
has a more delicate history and plays no role here.

For clarity, the displayed cyclotomic factor can be obtained directly,
without relying on a table of sign cases.  Suppose $|z|=1$ and
$z^a+z^b=1$.  The two unit vectors must be the two numbers
$e^{i\pi/3}$ and $e^{-i\pi/3}$, in either order.  If $w=z^d$, then
$w^A$ and $w^B$ are these primitive sixth roots.  A Bezout identity for
$A,B$ implies $w^6=1$, and $w^A$ of order six then implies that $w$
itself has order six.  Consequently $A,B$ have opposite residues
$1,5$ modulo six.  Conversely those residues make every root of
$X^{2d}-X^d+1$ a root of $F_{a,b}$.  These roots are simple: a common
zero of $F_{a,b}$ and its derivative on the unit circle would give
$a|z^a|=b|z^b|$, hence $a=b$, a contradiction.  Thus the formula for
$C_{a,b}$ accounts for the \emph{entire} cyclotomic factor.  The
remaining assertion, irreducibility of $H_{a,b}$, is the cited theorem.

\subsection{A conjugate-rotation invariant}

\begin{lemma}[The common divisor is intrinsic]
\label{rigid:lem:intrinsic-gcd}
Suppose $r^a+r^b=r^c+r^e=1$, where $0<r<1$,
$1\le a<b$, and $1\le c<e$.  Then
\[
\gcd(a,b)=\gcd(c,e).
\]
\end{lemma}

\begin{proof}
Set $d=\gcd(a,b)$.  Since $0<r<1$, it is not a root of unity, so
its minimal polynomial is $H_{a,b}$.  Theorem~\ref{rigid:thm:trinomial-input}
shows that this polynomial is a polynomial in $X^d$.  If $\zeta$ is a
primitive $d$th root of unity, $r\zeta$ is therefore a conjugate of $r$.
The second equation also holds at this conjugate:
\[
r^c\zeta^c+r^e\zeta^e=1.
\]
Taking absolute values gives equality in the triangle inequality,
because $r^c+r^e=1$.  Both summands must have the argument of their
positive real sum.  Thus $\zeta^c=\zeta^e=1$, proving
$d\mid\gcd(c,e)$.  Interchanging the pairs proves the reverse
divisibility.
\end{proof}

This step prevents a gap that would arise from normalizing the two
pairs independently.  The cyclotomic factor can have degree $2d$,
so a degree comparison before proving that the common divisors agree
would not suffice.

\subsection{The primitive coefficient comparison}

\begin{proof}[Proof of Theorem~\ref{rigid:thm:complement-collisions}]
First suppose an equal pair occurs.  Then $r^a=1/2$, so $r$ cannot
be an algebraic integer: otherwise the rational number $r^a=1/2$
would be an algebraic integer.  Any unequal pair would make $r$ a
root of the monic integer polynomial $F_{c,e}$, a contradiction.
There is only one equal pair because the powers of $r$ strictly
decrease.

Now suppose two distinct unequal pairs occur.  By
Lemma~\ref{rigid:lem:intrinsic-gcd} they have a common gcd $d$.
Replace $r$ by $r^d$ and divide all four exponents by $d$.
Both pairs are now primitive.  Their trinomials have the same
irreducible noncyclotomic factor, and their cyclotomic factors belong
to $\{1,\Phi_6\}$, where $\Phi_6=X^2-X+1$.
If the two cyclotomic factors agree, the trinomials agree, contrary
to distinctness.  After interchanging them if necessary, there must
therefore be a polynomial identity
\begin{equation}
(X^2-X+1)(X^n+X^m-1)=X^{n+2}+X^k-1,
\quad 1\le m<n,\quad 1\le k<n+2.
\label{rigid:eq:primitive-collision-factor}
\end{equation}
The coefficient of $X^{n+1}$ on the left is $-1$ unless
$m=n-1$, in which case it is zero.  The right-hand coefficient is
zero or one.  Hence $m=n-1$, and the identity reduces to
\[
X^{n+2}+X^{n-1}-X^2+X-1=X^{n+2}+X^k-1.
\]
For $n=2$ the remaining polynomial is $2X-X^2$, which is not a
monomial.  For $n>3$ its coefficient at $X^2$ is $-1$.
Thus $n=3$, which gives $m=2$ and $k=1$.
The common factor is $X^3+X^2-1$, and the two normalized pairs are
$(1,5)$ and $(2,3)$.  Conversely,
\[
X^5+X-1=(X^2-X+1)(X^3+X^2-1)
\]
shows that these two pairs do occur at $\omega$.
Finally, a third pair would have to be the same exceptional companion
of either existing pair, so no third pair occurs.
\end{proof}

\subsection{All signed pure-power unit equations}

The functional equations for polylogarithms use signs and inverses as
well as positive powers.  The preceding classification automatically
controls all such complementary arguments.

\begin{corollary}[Sharp signed-unit count]
\label{rigid:cor:signed-unit-count}
Let $\Gamma_r=\{\pm r^k:k\in\mathbb Z\}$ and
\[
\mathcal U(r)=\{x\in\Gamma_r:1-x\in\Gamma_r\}.
\]
Then $|\mathcal U(r)|$ is one of $0,3,6,12$.
For every unequal pair $(a,b)\in\mathcal C(r)$, its contribution is
the six-element set
\begin{equation}
\left\{r^a,r^b,r^{-a},r^{-b},-r^{b-a},-r^{a-b}\right\}.
\label{rigid:eq:six-unit-orbit}
\end{equation}
These sets are disjoint for different unordered pairs.
An equal pair contributes precisely $\{1/2,2,-1\}$.
The case $|\mathcal U(r)|=12$ occurs exactly when $r^d=\omega$
for some positive integer $d$.
\end{corollary}

\begin{proof}
The transformations $x\mapsto1-x$ and $x\mapsto1/x$ preserve
$\mathcal U(r)$.  For example,
$1-1/x=-(1-x)/x\in\Gamma_r$ whenever $x,1-x\in\Gamma_r$.
Their orbit consists of
\[
x,\quad1-x,\quad x^{-1},\quad(1-x)^{-1},
\quad x/(x-1),\quad(x-1)/x.
\]
Every real orbit has a representative in $(0,1)$: use $1/x$ if
$x>1$ and $1/(1-x)$ if $x<0$.  Such a representative and its
complement are $r^a,r^b$ with $a,b>0$.  The two members of a
nondegenerate real orbit that lie in $(0,1)$ are exactly $x$ and
$1-x$, so its unordered positive pair is unique.  Substituting
$x=r^a$, $1-x=r^b$ gives \eqref{rigid:eq:six-unit-orbit}.
All six entries are distinct unless $a=b$, when they reduce to
the stated three values.  The cardinalities now follow from
Theorem~\ref{rigid:thm:complement-collisions}.
\end{proof}

All four possibilities occur: $r=1/3$, $1/2$, the reciprocal golden
ratio, and $\omega$ give respectively $0,3,6,12$.
The conclusion classifies \emph{pure signed-power complements}.
It does not say that there are at most twelve useful polylogarithm
arguments: products of units $1-r^k$ can introduce many more arguments.

\begin{corollary}[Rational exponents]
\label{rigid:cor:rational-complements}
The same collision classification holds for positive rational
exponents, with the exceptional common scale $d$ allowed to be a
positive rational number.  In particular there are at most two
unordered positive rational pairs, and
\[
\#\{x\in\{\pm r^q:q\in\mathbb Q\}:1-x\in
                          \{\pm r^q:q\in\mathbb Q\}\}\le12.
\]
\end{corollary}

\begin{proof}
Given two pairs, choose a common denominator $N$ and apply
Theorem~\ref{rigid:thm:complement-collisions} to $r^{1/N}$.
A third rational pair could be included in the same denominator
and would contradict the integer theorem.  The orbit proof is
unchanged.
\end{proof}

\subsection{A provably complete bounded-degree search}

\begin{corollary}[Degree and an exhaustive exponent bound]
\label{rigid:cor:degree-search}
If $r^a+r^b=1$ with $1\le a<b$, then
\[
[\mathbb Q(r):\mathbb Q]=
\begin{cases}
b-2d,&(a/d,b/d)\equiv(1,5)\text{ or }(5,1)\pmod6,\\
b,&\text{otherwise},
\end{cases}
\qquad d=\gcd(a,b).
\]
Consequently all such bases of degree at most $D$ are found by
enumerating
\[
1\le a<b\le\left\lfloor\frac{5D}{3}\right\rfloor
\]
and retaining those whose displayed degree is at most $D$.
For primitive pairs it suffices to enumerate $b\le D+2$.
The only repeated bases in either enumeration are the classified
plastic pairs.  Equal pairs add the separate bases $2^{-1/a}$,
with $a\le D$.
\end{corollary}

\begin{proof}
The degree formula is Theorem~\ref{rigid:thm:trinomial-input}.  In the
exceptional case $b/d\ge5$, so $b-2d\ge3d$.  If the degree is
at most $D$, then $d\le D/3$ and
$b=(b-2d)+2d\le5D/3$.  The nonexceptional case is immediate.
For primitive pairs the possible degree loss is only two.
The equal-pair polynomial $2X^a-1$ has degree $a$ and is
irreducible: its reciprocal is, up to a nonzero scalar,
$X^a-2$, which is Eisenstein at $2$.
\end{proof}

This bound is sharp for infinitely many $D$: the pair $(d,5d)$
has degree $3d$ and largest exponent $5d$.
Unlike a search over bounded polynomial coefficients, the algorithm
has an explicit completeness theorem for its stated class.  No
numerical equality test is required.  A base is stored by the exact
irreducible polynomial $H_{a,b}$ together with its unique root in
$(0,1)$, and duplicate detection is exact polynomial comparison.

For illustration, the complete unequal-pair list through degree four is
\[
\begin{array}{ccl}
\text{degree}&\text{complement pair(s)}&\text{minimal polynomial}\\\hline
2&(1,2)&X^2+X-1\\
3&(1,3)&X^3+X-1\\
3&(1,5),(2,3)&X^3+X^2-1\\
4&(1,4)&X^4+X-1\\
4&(2,4)&X^4+X^2-1\\
4&(3,4)&X^4+X^3-1
\end{array}
\]
The degree-four entry with pair $(2,4)$ is exactly a square root of
the reciprocal golden ratio.  Its presence illustrates why normalization
of exponents is necessary when counting bases rather than primitive
seed types.

The exact enumeration in the accompanying certificate gives:
\[
\begin{array}{c|rrrrrrrrrrrr}
D&1&2&3&4&5&6&7&8&9&10&11&12\\\hline
\text{unequal-pair bases}&0&1&3&6&10&15&20&27&37&46&56&67\\
\text{primitive bases}&0&1&3&5&9&11&16&20&28&32&42&46\\
\text{all bases}&1&3&6&10&15&21&27&35&46&56&67&79
\end{array}
\]
Each row counts bases of degree at most $D$.  A primitive base has
an unequal complementary pair of gcd one; the intrinsic-gcd lemma
makes this definition independent of the choice of pair.
The last row also includes the $D$ equal-pair bases $2^{-1/a}$,
$1\le a\le D$.

\subsection{An exact counting formula}

The classification also counts bases without constructing or factoring
their minimal polynomials.  For $n\ge1$ define
\[
e(n)=\#\{a:1\le a<n,\ \gcd(a,n)=1,
            \ (a,n)\equiv(1,5)\text{ or }(5,1)\pmod6\}.
\]
In particular $e(n)=0$ unless $n\equiv1$ or $5\pmod6$,
and $e(1)=0$.

\begin{corollary}[Counting all bases by algebraic degree]
\label{rigid:cor:base-count}
Let $N(D)$ be the number of real $0<r<1$ with
$[\mathbb Q(r):\mathbb Q]\le D$ and $\mathcal C(r)\ne\varnothing$.
For every positive integer $D$,
\begin{align}
N(D)&=\frac{D(D+1)}2-\left\lfloor\frac D3\right\rfloor\notag\\
&\quad+\sum_{d=1}^{\lfloor D/3\rfloor}
\left\{e\!\left(\left\lfloor\frac Dd\right\rfloor+1\right)
       +e\!\left(\left\lfloor\frac Dd\right\rfloor+2\right)\right\}.
\label{rigid:eq:exact-base-count}
\end{align}
In particular,
\[
N(D)=\frac12D^2+O(D\log D),\qquad D\longrightarrow\infty.
\]
\end{corollary}

\begin{proof}
There are $D(D-1)/2$ unequal pairs with largest exponent at most
$D$, and their associated bases all have degree at most $D$.
A pair with largest exponent greater than $D$ can still have degree
at most $D$ only in the exceptional cyclotomic case.  Write that
pair $(dA,dB)$, where $\gcd(A,B)=1$ and the residues are opposite
$1,5$ modulo six.  The two inequalities are
\[
dB>D,\qquad d(B-2)\le D.
\]
Equivalently,
$B\in\{\lfloor D/d\rfloor+1,\lfloor D/d\rfloor+2\}$.
Since the exceptional normalized degree $B-2$ is at least three,
only $d\le D/3$ can occur.  This gives the sum in
\eqref{rigid:eq:exact-base-count}.
Every base counted twice is $\omega^{1/d}$, of degree $3d$;
there are exactly $\lfloor D/3\rfloor$ such duplicates.
Finally the $D$ equal-pair bases are distinct from all unequal-pair
bases and from each other.  This proves the exact formula.

For the asymptotic statement, use $0\le e(n)\le n$ to bound the sum
by
\[
\sum_{d\le D/3}\left(2\left\lfloor\frac Dd\right\rfloor+3\right)
\le2D\sum_{d\le D/3}\frac1d+D.
\]
The harmonic sum is $O(\log D)$, proving the claim.
\end{proof}

For an alternative exact implementation, when $n>1$ is coprime to
six, Mobius inversion gives
\begin{equation}
e(n)=\sum_{d\mid n}\mu(d)
       \left\lfloor\frac{n/d+1}{6}\right\rfloor.
\label{rigid:eq:exceptional-count-mobius}
\end{equation}
Indeed, after writing $a=dj$, the condition
$a\equiv-n\pmod6$ becomes $j\equiv-(n/d)\pmod6$ because every
unit modulo six is its own inverse.  For $t$ coprime to six,
the number of integers $1\le j<t$ in the residue class $-t$
is $\lfloor(t+1)/6\rfloor$.  Summing with $\mu(d)$ imposes
$\gcd(a,n)=1$.

% Citation keys: ZagierPolylogs1991, NakamuraShiraishi2025.

\section{Exact proofs and log-free ladders at the plastic base}
\label{rigid:sec:plastic-proofs}

The exceptional base in Theorem~\ref{rigid:thm:complement-collisions}
has two complementary pairs.  This additional structure produces
short proofs of two previously numerically supported formulas in
the manuscript, as well as log-free identities for $\zeta(2)$ and
$\zeta(3)$.

Throughout this section, $0<q<1$, $q^2+q^3=1$, and
$\lambda=\log q$.  The factorization in the preceding proof gives
\begin{equation}
1-q=q^5,\qquad1-q^2=q^3,\qquad
1-q^3=q^2,\qquad1-q^5=q.
\label{rigid:eq:plastic-complements}
\end{equation}
All logarithms used in the final identities are real.

\subsection{The dilogarithm formula: one Rogers pentagon}

Write
\[
R(x)=\operatorname{Li}_2(x)+\frac12\log x\log(1-x),
\qquad 0<x<1.
\]
The Rogers reflection and pentagon equations are
\begin{align}
R(x)+R(1-x)&=\zeta(2),\label{rigid:eq:plastic-R-reflect}\\
R(x)+R(y)&=R(xy)
 +R\!\left(\frac{x(1-y)}{1-xy}\right)
 +R\!\left(\frac{y(1-x)}{1-xy}\right).
\label{rigid:eq:plastic-R-pentagon}
\end{align}
They hold for $0<x,y<1$ and follow from the classical dilogarithm
functional equations; alternatively one can differentiate in $x$
and evaluate at $x=0$.  This real-domain formulation avoids branch
choices \cite{gaussian:Zagier}.

\begin{theorem}[The manuscript's plastic dilogarithm formula]
\label{rigid:thm:plastic-dilog}
\begin{equation}
\frac12\operatorname{Li}_2(q^2)-\operatorname{Li}_2(q^5)
=\lambda^2.
\label{rigid:eq:plastic-dilog-proved}
\end{equation}
\end{theorem}

\begin{proof}
Set $R_k=R(q^k)$.  At $(x,y)=(q,q^2)$ the three pentagon
arguments are
\[
xy=q^3,\qquad
\frac{x(1-y)}{1-xy}=q^2,\qquad
\frac{y(1-x)}{1-xy}=q^5.
\]
Hence $R_1+R_2=R_3+R_2+R_5$, or $R_1=R_3+R_5$.
Reflection gives $R_1+R_5=R_2+R_3=\zeta(2)$, so
$R_2=2R_5$.  Expanding the logarithmic corrections by
\eqref{rigid:eq:plastic-complements} gives
\[
\operatorname{Li}_2(q^2)+3\lambda^2
=2\operatorname{Li}_2(q^5)+5\lambda^2,
\]
as required.
\end{proof}

\subsection{The trilogarithm formula: one Kummer specialization}

For real nonzero $z\ne1$, define the single-valued trilogarithm
\[
P(z)=\operatorname{Re}\operatorname{Li}_3(z)
 -\log|z|\operatorname{Re}\operatorname{Li}_2(z)
 -\frac13\log^2|z|\log|1-z|,
\]
with the continuous value $P(1)=\zeta(3)$.
The inversion and duplication equations are
\[
P(z^{-1})=P(z),\qquad P(z)+P(-z)=\frac14P(z^2).
\]
Kummer's equation is \eqref{gaussian:eq:kummer-P3}, with $P=P_3$.
We use the form recorded by Zagier, Section~7, Example~3,
equation~(7) \cite{gaussian:Zagier}.

The single-valued Landen equation is
\[
P(x)+P(1-x)+P(1-x^{-1})=\zeta(3),\qquad0<x<1.
\]
Its ordinary real counterpart is
\begin{align}
\operatorname{Li}_3(x)+\operatorname{Li}_3(1-x)
 +\operatorname{Li}_3(1-x^{-1})
 &=\zeta(3)+\zeta(2)\log x\notag\\
 &\quad-\frac12\log^2x\log(1-x)+\frac16\log^3x.
\label{rigid:eq:plastic-ordinary-Landen}
\end{align}
One obtains it by differentiation using dilogarithm reflection and
Landen, with the integration constant determined at $x=1$.
See also Nakamura--Shiraishi \cite{gaussian:NakamuraShiraishi}.

\begin{theorem}[The manuscript's plastic trilogarithm formula]
\label{rigid:thm:plastic-trilog}
\begin{equation}
-\operatorname{Li}_3(q)+\frac32\operatorname{Li}_3(q^2)
 +\frac13\operatorname{Li}_3(q^3)-\operatorname{Li}_3(q^5)
=\lambda^3.
\label{rigid:eq:plastic-trilog-proved}
\end{equation}
\end{theorem}

\begin{proof}
Put $P_k=P(q^k)$.  At $(x,y)=(q,q^2)$, the nine arguments in
\eqref{gaussian:eq:kummer-P3}, in displayed order, are
\[
q,\quad q^2,\quad -q^{-1},\quad -q^4,\quad q^2,
\quad q^{-3},\quad q^3,\quad q^{-1},\quad q^{-5}.
\]
Using inversion and duplication, the equation becomes
\[
\mathsf K=-P_1+\frac92P_2+P_3-2P_4-P_5+\frac12P_8
=2\zeta(3).
\]
Landen at $x=q$, with $1-q=q^5$, similarly gives
\[
\mathsf A=P_1+P_5-P_4+\frac14P_8=\zeta(3).
\]
Consequently the exact row operation $(\mathsf K-2\mathsf A)/3$
proves
\begin{equation}
-P_1+\frac32P_2+\frac13P_3-P_5=0.
\label{rigid:eq:plastic-P3-proved}
\end{equation}

It remains to convert this to ordinary trilogarithms.  Denote the
left side of \eqref{rigid:eq:plastic-trilog-proved} by $D$, and put
\[
S=-\operatorname{Li}_2(q)+3\operatorname{Li}_2(q^2)
 +\operatorname{Li}_2(q^3)-5\operatorname{Li}_2(q^5).
\]
The two ordinary reflection equations and
Theorem~\ref{rigid:thm:plastic-dilog} give
\begin{align*}
\operatorname{Li}_2(q)+\operatorname{Li}_2(q^5)
 &=\zeta(2)-5\lambda^2,\\
\operatorname{Li}_2(q^2)+\operatorname{Li}_2(q^3)
 &=\zeta(2)-6\lambda^2,\\
\tfrac12\operatorname{Li}_2(q^2)-\operatorname{Li}_2(q^5)
 &=\lambda^2.
\end{align*}
Taking minus the first row, plus the second, plus four times the
third yields $S=3\lambda^2$.
The remaining logarithmic coefficient in the definition of $P$ is
\[
T=\{-1\cdot1^2\cdot5
 +(3/2)\cdot2^2\cdot3
 +(1/3)\cdot3^2\cdot2
 -1\cdot5^2\cdot1\}\lambda=-6\lambda.
\]
Thus the left side of \eqref{rigid:eq:plastic-P3-proved} is
\[
D-\lambda S-\frac{\lambda^2}{3}T
=D-3\lambda^3+2\lambda^3=D-\lambda^3.
\]
This proves \eqref{rigid:eq:plastic-trilog-proved}.
\end{proof}

The proof promotes the two displayed manuscript formulas from
numerical support to exact consequences of classical functional
equations.  It does not rely on pseudointegration at an isolated
algebraic base.

\subsection{Two log-free identities}

\begin{theorem}[Log-free plastic ladders]
\label{rigid:thm:plastic-logfree}
At the same base $q$, the following identities hold:
\begin{align}
&6\operatorname{Li}_2(q)-5\operatorname{Li}_2(q^2)
 -5\operatorname{Li}_2(q^3)+6\operatorname{Li}_2(q^5)
 =\zeta(2),\label{rigid:eq:plastic-logfree2}\\
&12\operatorname{Li}_3(q)-5\operatorname{Li}_3(q^2)
 -4\operatorname{Li}_3(q^3)-8\operatorname{Li}_3(q^4)
 +8\operatorname{Li}_3(q^5)+2\operatorname{Li}_3(q^8)
 =4\zeta(3).
\label{rigid:eq:plastic-logfree3}
\end{align}
\end{theorem}

\begin{proof}
Six times the first ordinary reflection equation minus five times
the second proves \eqref{rigid:eq:plastic-logfree2}.
For the trilogarithm, write $L_k=\operatorname{Li}_3(q^k)$.
Equation~\eqref{rigid:eq:plastic-ordinary-Landen} at $x=q$ and $x=q^2$,
with duplication to eliminate negative arguments, gives
\begin{align*}
A&=L_1+L_5-L_4+\frac14L_8
  =\zeta(3)+\zeta(2)\lambda-\frac73\lambda^3,\\
B&=\frac54L_2+L_3-L_1
  =\zeta(3)+2\zeta(2)\lambda-\frac{14}3\lambda^3.
\end{align*}
The logarithmic tails are proportional.  Therefore
$4(2A-B)=4\zeta(3)$, exactly the claimed six-argument identity.
\end{proof}

These formulas are explicit additional identities for the package.
Their derivations and the classical origins of the functional
equations are given in full; no global priority claim is made.
They also give exponentially convergent algebraic series, for example
\begin{equation}
\zeta(3)=\sum_{n=1}^{\infty}
\frac{3q^n-\tfrac54q^{2n}-q^{3n}-2q^{4n}
                 +2q^{5n}+\tfrac12q^{8n}}{n^3}.
\label{rigid:eq:plastic-zeta3-series}
\end{equation}
Absolute convergence justifies expansion of the polylogarithms.
If this series is truncated after $N$ terms, its absolute tail is
at most
\[
\frac{39}{4}\,
\frac{q^{N+1}}{(N+1)^3(1-q)},
\]
obtained by summing the absolute coefficients and using
$q^{kn}\le q^n$ for $k\ge1$.  The constant $39/4$ is deliberately
simple; the sharper sum of the six individual geometric tails can
also be used.


\section{The modified polylogarithm}
Zagier's one-valued modified polylogarithm (his eq.~(33)) is
\begin{equation}\label{bloch:eq:Pm}
  P_m(x)=\Real_m\!\Bigl(\sum_{j=0}^{m-1} \frac{2^j B_j}{j!}\,(\log|x|)^j\,\Li_{m-j}(x)\Bigr),\qquad
  \Real_m=\begin{cases}\Real,&m\text{ odd},\\ \operatorname{Im},&m\text{ even},\end{cases}
\end{equation}
with $B_j$ the Bernoulli numbers and $\Li_1(x)=-\log(1-x)$. It is continuous on the Riemann sphere away from
$\{0,1,\infty\}$, satisfies $P_m(\bar x)=(-1)^{m-1}P_m(x)$, and $P_m(1)=\zeta(m)$ for odd $m$, $0$ for even $m$;
$P_2(\mathrm i)$ is Catalan's constant. We verified our implementation against Zagier's worked examples
$P_5(-\tfrac18)-126P_5(\tfrac12)-162P_5(-\tfrac12)=\tfrac{403}{16}\zeta(5)$ (his eq.~(825)) and the long
thirteen-term relation (825)/(833), to $60$ digits.

\section{Ladders are Bloch-group elements}
The Bloch-group viewpoint records the obstruction to simplifying a
dilogarithm combination. At weight two the boundary sends
$[x]\mapsto x\wedge(1-x)$; a vanishing boundary makes a combination a
candidate for the Bloch group, modulo the five-term relation. Higher weights
require further relations and compatibility conditions: a weight-one unit
relation alone does not prove that its pseudointegrations are higher
Bloch-group elements. Modified polylogarithms remove explicit logarithmic
tails and make the following ladder combinations transparent:

\begin{align}
  \text{golden}:&\quad 15\,P_3(\phi^{-2})=12\,\zeta(3),\label{bloch:eq:goldenP3}\\
  \text{supergolden}:&\quad -60\,P_3(\rho)+54\,P_3(\rho^2)-32\,P_3(\rho^3)-3\,P_3(\rho^6)=-36\,\zeta(3),
  \label{bloch:eq:sgP3}
\end{align}
both exact. Compare~\eqref{bloch:eq:goldenP3} with the ``raw'' form $15\Li_3(\phi^{-2})=10\ln^3\phi-2\pi^2\ln\phi+12\zeta(3)$:
the modified trilogarithm removes the explicit logarithmic tail.
At odd weight, comparable modified combinations require the appropriate
lower-weight identities. At even weight the definition takes an imaginary
part, and every term is real on $0<x<1$, so $P_{2j}(x)=0$. The ordinary
even-weight golden formulas retain their content; under this definition
they cannot become nonzero even-zeta multiples. This distinction concerns
the chosen modification, not a universal bound on regulator directions.

\section{The bases as number fields}
\begin{center}\small
\begin{tabular}{lllll}
\toprule
base $\rho$ & minimal polynomial of $1/\rho$ & field $F$ & signature & disc \\
\midrule
golden $\phi^{-1}$ & $x^2-x-1$ & $\mathbb Q(\sqrt5)$ & $(2,0)$ totally real & $5$ \\
minimal Pisot $\omega$ & $x^3-x-1$ & $\mathbb Q(\theta)$ & $(1,1)$ & $-23$ \\
supergolden recip.\ & $x^3-x^2-1$ & $\mathbb Q(\sigma)$ & $(1,1)$ & $-31$ \\
\bottomrule
\end{tabular}
\end{center}
The middle row is Zagier's own running example (\S4--5, \S9C): the field $\mathbb Q(\theta)$, $\theta^3-\theta-1=0$,
whose ladder he calls ``particularly spectacular,'' and whose defining unit relations
$1-\theta=-\theta^{-4},\,1+\theta=\theta^3,\dots$ are precisely the seeds we rediscovered. Our supergolden base
gives a \emph{different} complex cubic field, discriminant $-31$.

\section{Regulator parity and the limits of ladder counting}
For $m\ge2$, Borel's rank theorem gives rank $r_1+r_2$ for
$K_{2m-1}(F)$ when $m$ is odd, and $r_2$ when $m$ is even.
The modified polylogarithm uses its real part at odd weight and imaginary
part at even weight. In particular $P_{2m}$ vanishes on real arguments,
even when the number field has nonzero regulator directions at its complex
embeddings. This explains why a real-argument search can miss arithmetic
content visible at a complex embedding.

Cyclotomic unit searches compute
$\operatorname{Norm}\Phi_k(\rho)=\operatorname{Res}(p,\Phi_k)$ and select
the unit cases. For a fixed finite seed set, imposing successive regulator
constraints is a finite linear-algebra problem. A dimension count can give
a lower bound for a kernel or exhaust that chosen construction. It does
not prove there are no additional seeds, no dependencies among constraints,
and no other functional identities at higher weight.

The later project source identifies minimal-Pisot ladders through
weight six. The published literature already includes the 1987
Abouzahra--Lewin--Xiao paper~\cite{golden:AbouzahraLewinXiao1987},
which distinguishes analytic low-weight results from numerical higher
ladders, and Gangl's proved weight-six and weight-seven relations for
single-valued modified polylogarithms~\cite{golden:Gangl2013}.
The modified and ordinary polylogarithms must be distinguished when
comparing these statements, particularly at real arguments and even
weight. Cohen--Lewin--Zagier's published numerical sixteenth-order
Lehmer ladder~\cite{CohenLewinZagier} supplies another substantial example.
These results supersede earlier bounded negative searches. The project
source's stronger assertion that the supergolden field has no weight-five
real ladder is not justified by exhausting its selected seed set, and
is omitted. A maximal-weight theorem needs completeness for the chosen
class of higher-weight relations; completeness of multiplicative seeds
alone does not establish it.

\section{Attributed dilogarithm and inverse-binomial examples}\label{sec:dilog-examples}
The following examples are attributed to Hakimoglu-Brown~\cite{HakimogluBrown}. Their displayed numerical checks in the source are historical.

All to $\ge38$ digits (\code{tools/scratch/dilog-inverse-binomial/}); a representative selection:
\begin{itemize}
\item \textbf{Sun's open conjecture} (eq.\ (24)), proved in the paper via the $\mathbb{Q}(\sqrt3)$
  evaluation (25):
  \[ \sum_{n\ge0}\frac{\binom{2n}{n}}{(2n+1)8^{n}}\Bigl(\sum_{0\le k<n}\frac{(-1)^k}{2k+1}-\frac{(-1)^n}{2n+1}\Bigr)=-\frac{\sqrt2}{16}\pi^2. \]
\item \textbf{Theorem 2} ($\mathbb{Q}(\sqrt3)$):
  $8\Li_2(2-\sqrt3)-\Li_2(4\sqrt3-7)=\tfrac{5\pi^2}{12}+\ln(2-\sqrt3)\ln(2+\sqrt3)$,
  and $3\Li_2(7-4\sqrt3)=6\Li_2(2-\sqrt3)+6\Li_2(\sqrt3-2)$.
\item \textbf{Lima's identity} ($\mathbb{Q}(\sqrt2)$, eq.\ (26)):
  $\Li_2(\tfrac1{\sqrt2})-\Li_2(1-\sqrt2)=\tfrac{\pi^2}6-\tfrac{\ln^2 2}8+\tfrac12\ln2\,\ln(\sqrt2-1)$.
\item \textbf{Zagier's four-term} (eq.\ (27)): $3L(-\tfrac16)-L(\tfrac18)+L(\tfrac19)+L(\tfrac1{28})=-\tfrac{\pi^2}{12}$,
  with the (unnormalized) Rogers $L(x)=\Li_2(x)+\tfrac12\ln|x|\ln|1-x|$; and its compressed three-term form.
\item the rational-argument $\tfrac{\pi^2}{18}+2\ln2\ln3-2\ln^2 2-\tfrac23\ln^2 3=\Li_2(\tfrac14)+\tfrac13\Li_2(\tfrac19)$,
  the inverse-tangent/$\Cl_2(\tfrac\pi3)$ relation (30), and $\Li_2(\tfrac43)+\Li_2(4)=\tfrac12(\pi^2-\ln^2 3)$ (real parts).
\item the \textbf{engine}: $\,{}_4F_3(\{1,1,1,\tfrac32\};\{\tfrac52,\tfrac52,\tfrac52\};1)=\tfrac{27}2\bigl(7\zeta(3)+(3-2G)\pi-12\bigr)$
  (D'Aurizio) and the cube-binomial $_4F_3$ of eq.\ (1).
\item the paper's \textbf{claimed-original elliptic $_4F_3(1)$ evaluations} (eqs.\ 32, 33, 35):
  \begin{align*}
   \Cl_2(\tfrac\pi3)&=\tfrac{\sqrt3}2\,{}_4F_3\bigl(\{1,\tfrac12,\tfrac12,\tfrac12\};\{\tfrac32,\tfrac56,\tfrac76\};1\bigr),\\
   {}_4F_3\bigl(\{1,\tfrac12,\tfrac23,\tfrac13\};\{\tfrac32,\tfrac54,\tfrac34\};1\bigr)&=\tfrac{3\sqrt3}4\bigl(2\ln(4+3\sqrt2)+\ln(\tfrac72-2\sqrt3)\bigr),\\
   \ln^2\!\Bigl(\tfrac{5+\sqrt{13}}{2\sqrt3}\Bigr)-\tfrac{\pi^2}{12}&=\tfrac1{180}\,{}_4F_3\bigl(\{1,1,1,\tfrac32\};\{\tfrac76,\tfrac{11}6,2\};-\tfrac1{324}\bigr).
  \end{align*}
\end{itemize}
\textbf{All check out.} One caution worth recording: equation (27) uses the \emph{unnormalized} Rogers
$L$ (so the right side is $-\tfrac{\pi^2}{12}$, not a pure rational); with the $\tfrac6{\pi^2}$-normalized
$L$ the identity reads $3L(-\tfrac16)-L(\tfrac18)+L(\tfrac19)+L(\tfrac1{28})=-\tfrac12$.

\paragraph{Connection.} The paper evaluates inverse-binomial series with $\binom{2n}{n}^{\!2}$ and
$\binom{3n}{n}$ via beta integrals; this is the reciprocal-binomial sibling of the central-binomial
ladder $\sum\binom{2n}{n}4^{-n}(2n{+}1)^{-p}$ studied in \code{harvested-central-binomial-arcsin.wl}.
The verified identities are stored (with attribution) in
\code{src/PolyLog/identities/harvested-dilogarithm-ladders-2506.10206.wl}.
